% =============================================================================
% Das Brücken-Zertifikat -- ein zeilengenaues Prüfinstrument für
% theorieübergreifende Transportargumente, mit vier Anwendungen auf die
% abc-Literatur.
% Entwurf v0.2 (Deutsch) -- 1:1-Fassung zur führenden englischen Version.
% Projekt: abc_forensik/paper_bridge_certificate  (.RESEARCH/.LAB/.HCT)
% Teil 3 der abc-Gap-Serie (Zählung nach Serienregister, Entscheid E2=a vom
% 2026-08-13). KEIN Beweispapier und KEIN Fehler-Nachweis.
% Siehe die Geltungsbereichs-Box nach dem Abstract.
%
% v0.2 (2026-08-16): übernimmt die gehärtete sechste Vertragskomponente des
% inzwischen veröffentlichten Methodenpapiers (Teil 2, v0.2.4). Zwei Verdikte
% aus v0.1 werden umcodiert -- Fall A Pflicht 6 (eingelöst -> liefert nicht)
% und Fall C Pflicht 6 (eingelöst -> teilweise) -- und die Umwertung wird im
% unnummerierten Abschnitt "Änderungen in Version 0.2" nach dem
% Inhaltsverzeichnis offen ausgewiesen. Ferner: Split R6a/R6b, Versionsetikett
% und Namespace-Hinweis für das Instrument, R8 als getriggerte Erweiterung (E8)
% klargestellt, Fall C umbenannt in "vergleichende Diskriminierungsanwendung",
% LANA-Mapping primär auf Pflicht 6 gebucht, Teil 2 als veröffentlicht zitiert.
% v0.1 bleibt auf Zenodo unverändert erhalten.
% =============================================================================

\documentclass[11pt,a4paper]{article}

\usepackage[utf8]{inputenc}
\usepackage{lmodern}   % skalierbare Fonts -- nötig für microtype-Expansion
\usepackage{cmap}
\usepackage[T1]{fontenc}
\usepackage[ngerman]{babel}
\usepackage[a4paper,margin=2.7cm]{geometry}
\usepackage{amsmath}
\usepackage{amsthm}
\usepackage{amssymb}
\usepackage{array}
\usepackage{booktabs}
\usepackage{microtype}
\usepackage{xurl}
\usepackage[colorlinks=true,linkcolor=black,citecolor=black,urlcolor=blue]{hyperref}

\input{glyphtounicode}
\pdfgentounicode=1

\hyphenation{Teich-mül-ler in-ter-uni-ver-sell In-de-ter-mi-na-cy
  mul-ti-ra-di-al Pro-zes-sions-nor-mie-rung Nor-mie-rungs-wahl
  Arith-me-ti-coid Arith-me-ti-coide Zer-ti-fi-kat}

\emergencystretch=1.5em

\newcolumntype{P}[1]{>{\raggedright\arraybackslash}p{#1}}

\hypersetup{
  pdfencoding=unicode,
  psdextra,
  pdftitle={Das Brücken-Zertifikat},
  pdfauthor={Lukas Geiger},
  pdfsubject={Ein zeilengenaues Prüfinstrument für theorieübergreifende Transportargumente}
}

% Sprachspezifisches Zielnamen-Präfix, damit das zweisprachige Kombi-PDF die
% Anker beider Sprachfassungen auseinanderhält.
\makeatletter
\renewcommand*{\HyperDestNameFilter}[1]{bc-de-#1}
\makeatother

% --- Umgebungen ---------------------------------------------------------------
\theoremstyle{definition}
\newtheorem{obligation}{Pflicht}
\newtheorem*{workinglemma}{Arbeitslemma (Brückenvertrag)}
\newtheorem*{recipe}{Anwendungsrezept}

% --- Notation -----------------------------------------------------------------
\newcommand{\Brxi}{\mathrm{Br}_{\xi}}
\newcommand{\LogVol}{\operatorname{LogVol}}
\newcommand{\Vol}{\operatorname{Vol}}
\newcommand{\ordd}{\operatorname{ord}}
\newcommand{\Thet}{\widetilde{\Theta}}

% Jede Quellenangabe ist an einen benannten Extrakt gebunden, weil die hier
% verwendeten Extrakte unabhängig durchnummeriert sind; eine bloße Zeilennummer
% wäre mehrdeutig.
\newcommand{\anchor}[2]{(#1, Extrakt Z.~#2)}

\title{Das Brücken-Zertifikat:\\
Ein zeilengenaues Prüfinstrument für theorieübergreifende\\
Transportargumente, mit vier Anwendungen\\
auf die $abc$-Literatur}
\author{Lukas Geiger\\[2pt]
\small Unabhängiger Forscher, Bernau, Deutschland\\
\small ORCID: 0009-0005-7296-1534}
\date{Oktober 2026 --- Version 0.3 (korrigierter Entwurf)}

\begin{document}
\maketitle

\begin{abstract}
\noindent
Eine wiederkehrende Argumentform der arithmetischen Geometrie trägt einen
numerischen Readout aus einer Konstruktion heraus in einen Container, den sie
mit einer zweiten teilt, und liest daran eine Ungleichung ab. Streitfälle über
solche Schritte haben sich als ungewöhnlich widerstandsfähig gegen die üblichen
Instrumente der mathematischen Begutachtung erwiesen: Die Beteiligten sind sich
über nahezu jede Formel einig und uneins darüber, ob eine bestimmte Passage
überhaupt ein Argument \emph{ist}. Diese Arbeit schlägt für genau diese Lage ein
Instrument vor und berichtet, was es in vier Anwendungen zurückgegeben hat. Das
Instrument --- ein \emph{Brücken-Zertifikat} --- legt sieben Pflichten des
vorgeschlagenen Zertifikats für Übergänge von einer qualitativen Verknüpfung
zu einer numerischen Membership-Aussage fest, ordnet sie als neunzeiliges Ledger an, das
zeilenweise gegen den Wortlaut eines Primärtextes gefüllt wird, und bindet die
Befunderfassung an Regeln: Verdikte betreffen die Zertifikatsform, nie die
mathematische Wahrheit; jeder Befund trägt ein wörtliches Zitat mit Zeilenanker;
Lücken werden typisiert statt graduiert; jeder Rückgriff auf Standardliteratur
wird in ein Import-Ledger eingetragen; und jede Aussage der Form \glqq das steht nicht
im Text\grqq{} wird zusammen mit den verwendeten Suchbegriffen berichtet. Vier
Anwendungen werden vorgeführt, und sie leisten Verschiedenes. Am Beweis eines
Korollars der interuniversellen Teichmüller-Theorie \emph{lokalisiert} das Ledger
eine nicht eingelöste Pflicht: Die entscheidende Membership folge, so der Text,
\glqq formal\grqq{} aus einem Paar deklarierter Eigenschaften, statt abgeleitet zu werden.
Am Folgeband, den derselbe Text als seinen Quantifizierungsort benennt,
\emph{schließt} das Ledger diesen Ort \emph{aus}: Die dortige Rechnung ist
explizit und an jedem von ihr benannten Knoten geschlossen --- nachgerechnet hat
diese Arbeit sie nicht ---, verbraucht die strittige Ungleichung aber als Eingabe,
statt sie zu erzeugen. An einem unabhängigen Vergleichsprogramm in einer anderen
Theorie
dient das Ledger als \emph{vergleichende Diskriminierungsanwendung} und liefert
eine anders geformte Lücke an anderer Adresse. An zwei nachgelagerten Anwendungen protokolliert es
\emph{Vererbung}: Der strittige Schritt wird in einem Satz importiert und nie
diskutiert. Im selben Zeitraum hat ein internationales Formalisierungsprojekt, das
diese Arbeit nicht zitiert, an derselben Passage eine
Kompatibilitätsaussage isoliert; wir ordnen ihre Objekte den vorgeschlagenen
Rollen der Zertifikatspflichten zu. Dies beweist keine Implikation von der
Kompatibilitätsgleichung des Berichts zur typisierten Ordnungsbedingung
des Zertifikats. Nichts in dieser Arbeit ist die Behauptung, ein geprüftes
Argument sei falsch, und nichts darin ist eine Aussage über die
$abc$-Vermutung.
\end{abstract}

% -----------------------------------------------------------------------------
% GELTUNGSBEREICH UND CLAIM-GRENZE
% -----------------------------------------------------------------------------
\begin{center}
\fbox{%
\parbox{\dimexpr\textwidth-2\fboxsep-2\fboxrule\relax}{%
\medskip
\textbf{Geltungsbereich und Claim-Grenze.}\par\medskip
\begin{itemize}\setlength{\itemsep}{2pt}
\item Diese Arbeit beweist \emph{nichts} über die $abc$-Vermutung. Sie enthält
  keinen Beweis und keine Widerlegung, und keine Aussage darüber, ob die
  Vermutung mit den behandelten Methoden beweisbar ist.
\item Sie enthält \emph{keinen Fehler-Nachweis}. Eine Ledger-Zeile mit dem
  Verdikt \glqq liefert nicht\grqq{} hält fest, dass eine benannte Pflicht \emph{im
  geprüften Wortlaut} nicht eingelöst wird. Das ist eine Aussage über die Form
  eines Zertifikats, nicht über die Wahrheit eines Satzes. \glqq Liefert nicht\grqq{} ist
  nicht \glqq widerlegt\grqq.
\item Sie \emph{verifiziert keine Schranke}. Keine Beweiszeile wurde
  nachvollzogen, keine Abschätzung nachgerechnet, keine Konstante geprüft. Wo
  berichtet wird, dass in einer Quelle gerechnet wird, ist das ein Befund
  darüber, \emph{dass} und \emph{wo} gerechnet wird, keine Bestätigung, dass es
  richtig ist.
\item Sie ergreift in keinem bestehenden Streit über die geprüften Argumente
  Partei. Bekannte Einwände gehen hier ausschließlich als Definition einer
  Negativkontrolle ein, nie als Beleg.
\item \textbf{Kein Autoritätsargument.} Weder die Selbstauskunft eines Autors,
  ein Schritt genüge, noch die Auskunft eines Kritikers, er genüge nicht, zählt
  hier als Beleg. In den geprüften Texten gefundene Selbsteinschätzungen sind
  Prüfgegenstand, nicht Stütze.
\item \textbf{Textbasis} sind \texttt{pdftotext}-Extrakte der Quell-PDFs. Alle
  Zeilenangaben (\glqq Extrakt Z.~NNNN\grqq) sind Auffindhilfen in diese Extrakte und
  \emph{keine} Seitenzählung der publizierten Dokumente. Anders als bei einer
  rein prosabezogenen Prüfung hängt hier \emph{ein} Verdikt aus
  Abschnitt~\ref{sec:caseA} an Symbolen, die bei der Extraktion verloren gingen;
  die Rekonstruktion wird in Abschnitt~\ref{sec:reconstruction-A} offengelegt,
  und die zugehörige Regel ist Bestandteil des Instruments
  (Abschnitt~\ref{sec:reconstruction}).
\item Die bibliographischen Einträge sind vorläufig und tragen ihre Provenienz;
  drei von ihnen sind bewusst auf das Belegte reduziert. Siehe
  Abschnitt~\ref{sec:limitations}, Punkt~\ref{lim:biblio}.
\end{itemize}
\medskip
}}
\end{center}

\clearpage
\tableofcontents
\clearpage

% =============================================================================
% ÄNDERUNGEN IN VERSION 0.2 -- bewusst unnummeriert: die Abschnittsnummerierung
% der veröffentlichten Version 0.1 darf sich nicht verschieben.
% =============================================================================
\phantomsection
\section*{Änderungen in Version 0.3}
\addcontentsline{toc}{section}{Änderungen in Version 0.3}
\label{sec:changes03}

Diese Korrekturversion trennt das prospektive Rezept von den historischen
Teilanwendungen, korrigiert den Fall-C-Verweis der Synopse auf ATS IV v2
\S\,6.10/Theorem~6.10.1 und stellt die Originalnotation der LANA-Gleichung
$\eta^q=\eta_S^{\mathrm{anab}}$ (9-1) wieder her.
Der LANA-Abgleich ist eine vorgeschlagene Rollenzuordnung: Weder ein
Volumenquotient noch die Gleichheit abstrakter Abbildungen beweist für sich
die typisierte Halbgeradenbedingung des Zertifikats. Die erforderliche
Objekt-, Auswertungs- und Schrankenbrücke bleibt offen.
Die sieben Pflichten definieren den vorgeschlagenen Formvertrag; ihre
universelle Notwendigkeit für jeden möglichen mathematischen Beweis ist
nicht gezeigt. Fall D betrifft nachgelagerte Zhou-Importe, keine vollständige
C7-Positivkontrolle oder nachgewiesene zyklische Abhängigkeit. Fall C beweist
keinen mathematischen Widerspruch in ATS.
Die schon vorbereiteten nummerierten Tabellen mit Beschriftungen und zwei
Synopsen bleiben erhalten. Englischer und deutscher Text enthalten dieselben
Korrekturen. Kein neues mathematisches Resultat und keine unabhängige
Methodenvalidierung werden beansprucht.

\medskip

\section*{Änderungen in Version 0.2}
\addcontentsline{toc}{section}{Änderungen in Version 0.2}
\label{sec:changes02}

Version 0.1 dieser Arbeit erschien, bevor das Methodenpapier der
Serie~\cite{GeigerBCM2026} fertiggestellt war. Jenes Papier hat die sechste
Vertragskomponente gehärtet: Eine gemeinsame numerische Auswertung im Container
löst sie nicht mehr für sich allein ein; die Kompatibilitätsidentitäten, die die
Behauptung braucht --- insbesondere $\nu_C \circ i_q = \nu_q$ auf dem
ausgewiesenen Quellbereich ---, sind zwingend, und eine stipulierte Invarianz ist
kein Beweis einer solchen Identität. Version 0.2 übernimmt diese Lesart, und mit
ihr ändern sich zwei Verdikte, die Version 0.1 protokolliert hatte. Eine
veröffentlichte Codierung darf nicht stillschweigend revidiert werden, deshalb
werden beide hier benannt.

\begin{itemize}
\item \textbf{Fall~A, Pflicht~6: eingelöst $\rightarrow$ liefert nicht.} Version
  0.1 protokollierte Pflicht~6 als \glqq über ihre erste Alternative\grqq{} eingelöst, weil
  eine numerische Auswertung im Container existiert, und hielt fest, \glqq die
  Fehlstelle sitzt bei Pflicht~7, nicht bei Pflicht~6\grqq. Unter der
  Kompatibilitätslesart ergibt dieselbe Evidenz die umgekehrte Codierung --- und
  die Evidenz stammt aus der Arbeit selbst: Der Block \glqq gesucht, nicht gefunden\grqq
  in Abschnitt~\ref{sec:caseA} hielt bereits fest, dass die Identität
  $\nu_C \circ i_q = \nu_q$ über die Container-Identifikation hinweg nicht im
  Text steht und dass das eine vorhandene Kompatibilitätslemma auf einer anderen
  Achse liegt. Für diese Änderung wurde keine Quelle neu geprüft und keine neue
  Evidenz erhoben; geändert hat sich die Regel, gegen die die vorhandene Evidenz
  gelesen wird.
\item \textbf{Fall~C, Pflicht~6: eingelöst $\rightarrow$ teilweise.} Version 0.1
  protokollierte Pflicht~6 als eingelöst, wieder über die erste Alternative. Die
  Auswertung dort ist tatsächlich numerisch, aber die Kompatibilität hält nur
  zwischen den Summen und wird stellenweise von der Quelle selbst dementiert ---
  in einer Passage, die Version 0.1 einen Absatz weiter bereits zitierte.
\end{itemize}

Drei Folgen werden im Text mitgeführt statt auf diese Notiz beschränkt: Pflicht~6
erhält eine eigene Ledger-Zeile (der Split R6a\,/\,R6b in
Abschnitt~\ref{sec:splits}); das Anwendungsrezept in Abschnitt~\ref{sec:recipe}
wird in Schritt~4 korrigiert, damit eine Nachnutzung nicht die abgelöste Lesart
reproduziert; und die zusammenfassenden Aussagen, die den Fall-A-Befund mit einer
eingelösten Pflicht~6 flankierten --- in der Ergebnistabelle, in der Rollentabelle
von Abschnitt~\ref{sec:convergence} und im Schluss ---, werden nachgezogen.

\paragraph{Was sich nicht geändert hat.} Der zentrale Befund bleibt unberührt:
Pflicht~7 ist am geprüften Substep nicht eingelöst, und die Membership soll
\glqq formal\grqq{} aus einem Eigenschaftspaar folgen. Das Abstract von Version 0.1 trug den
Pflicht-6-Anspruch nicht und bleibt der Sache nach unverändert. Die Fälle~B
und~D behalten ihre Verdikte, und kein Fall wurde neu auditiert.

\paragraph{Die übrigen Änderungen sind terminologisch} und gleichen diese Arbeit
an das fertiggestellte Methodenpapier an. Das Instrument erhält ein
Versionsetikett, und seine präfixlosen Zeilenlabels werden eingeordnet
(Abschnitt~\ref{sec:ledger}). Zeile R8 wird als getriggerte Erweiterung statt als
Kernzeile ausgewiesen, was den Nenner korrigiert, ohne ihre Codierung zu ändern.
Und Fall~C wird von \glqq Vergleichskontrolle\grqq{} in \emph{vergleichende
Diskriminierungsanwendung} umbenannt, den Begriff von Teil~2 --- eine reine
Umbenennung, denn diese Arbeit hielt bereits fest, dass eine echte
Positivkontrolle nicht unter den geprüften Fällen war.

\clearpage

% =============================================================================
\section{Einleitung}
\label{sec:intro}
% =============================================================================

\subsection{Eine Meinungsverschiedenheit, die gewöhnliche Begutachtung nicht auflöst}

Es gibt eine Klasse mathematischer Meinungsverschiedenheiten, bei der der übliche
Reparaturmechanismus versagt. Es ist nicht so, dass eine Seite sich verrechnet
hätte und die andere den Fehler fände. Es ist nicht so, dass ein Lemma falsch
wäre und ein Gegenbeispiel die Sache entschiede. Es ist so, dass eine Textpassage
von ihrem Autor als Argument und von einem Leser als Umformulierung gelesen wird
und dass beide Lesarten den Kontakt mit den Formeln überstehen --- weil die
Formeln nicht strittig sind.

Die Passage im Zentrum dieser Arbeit hat diesen Charakter. Im dritten Band der
interuniversellen Teichmüller-Theorie erreicht der Beweis eines Korollars einen
Punkt, an dem zwei Größen --- ein numerischer Readout an einem Objekt und eine
durch ein zweites Objekt bestimmte reelle Halbgerade --- in eine einzige
Auswertung gebracht werden, und dann wird die Enthaltensein-Beziehung der ersten
in der zweiten behauptet. Um diese Enthaltensein-Beziehung herum ist eine ganze
arithmetische Konsequenz organisiert. Ob die Passage sie herstellt oder
voraussetzt, wird seit Jahren ohne Konvergenz diskutiert.

Gewöhnliche Begutachtungsprosa ist hier aus strukturellen Gründen ein schlechtes
Instrument. Sie beantwortet die Frage \glqq ist das korrekt?\grqq. Wenn der Streit darum
geht, ob eine gegebene Satzfolge überhaupt eine Ableitung ist, hat diese Frage
keine stabile Bedeutung: Jede Seite beantwortet sie über einen anderen
Gegenstand. Schlimmer noch, Prosabegutachtung erlaubt --- ja lädt dazu ein --- einen
Zug, der den Streit prinzipiell unauflösbar macht, nämlich den Rückgriff auf
Autorität. \glqq Der Autor sagt, die Konstruktion genüge\grqq{} und \glqq die Kritiker sagen, sie
genüge nicht\grqq{} sind beides Aussagen über Personen. Sobald sie eintreten, ist der
Text nicht mehr der Gegenstand.

\subsection{Was ein Prüfinstrument stattdessen leistet}

Die hier verfolgte Alternative besteht darin, nicht mehr zu fragen, ob das
Argument korrekt ist, sondern eine Frage zu stellen, die ein Text aus sich heraus
beantworten kann: \emph{Welche Pflichten muss eine Passage dieser Art einlösen,
und welche davon löst dieser konkrete Wortlaut ein?}

Diese Umformulierung erkauft vier Dinge.

\begin{enumerate}
\item \textbf{Sie ist am Text entscheidbar.} Ob ein Dokument ein Lemma enthält,
  das eine bestimmte Implikation behauptet, ist eine Frage über das Dokument.
  Zwei sorgfältige Leser können darüber streiten, ob ein Beweis überzeugt; sie
  werden selten darüber streiten, ob ein Satz vorhanden ist.
\item \textbf{Sie lokalisiert.} Ein Verdikt hängt an einer nummerierten Pflicht
  und über sie an einem Zeilenbereich eines bestimmten Textes. Aus \glqq das Argument
  hat eine Lücke\grqq{} wird \glqq Pflicht~7 ist an dieser Passage nicht eingelöst, und hier
  ist der Satz, der an ihrer Stelle steht\grqq.
\item \textbf{Sie ist symmetrisch.} Dasselbe Instrument wird unverändert auf
  einen Vorschlag und auf seinen Konkurrenten angewandt. Eine Methode, die für
  einen Text \glqq fällt durch\grqq{} liefert und für einen anderen kein Verfahren hat, ist
  kein Instrument, sondern ein Urteil auf der Suche nach einem Gegenstand.
\item \textbf{Sie begrenzt ihre eigenen Behauptungen.} Weil die Frage die Form
  betrifft, kann die Antwort nicht stillschweigend zu einer Aussage über Wahrheit
  anwachsen. Eine nicht eingelöste Pflicht ist ein fehlender
  Zertifikatsbestandteil, keine Widerlegung --- und das Instrument ist so gebaut,
  dass diese Unterscheidung in der Berichterstattung nicht verlorengehen kann.
\end{enumerate}

Der Preis ist ebenso klar und wird hier einmal genannt, weil alles in dieser
Arbeit ihn erbt: Ein Instrument dieser Art kann niemals feststellen, dass ein
Argument stichhaltig ist. Alle Pflichten können eingelöst sein und das Argument
kann dennoch aus einem Grund falsch sein, nach dem das Instrument nicht sucht. Es
kann nur feststellen, dass eine bestimmte, benannte Anforderung von einem
bestimmten, zitierten Wortlaut erfüllt wird oder nicht.

\subsection{Das Instrument in einem Absatz}

Ein \emph{Brücken-Zertifikat}, geschrieben $\Brxi$, ist eine Liste von sieben
Pflichten, die an einen Argumentationsschritt bestimmter Gestalt geknüpft sind:
an einen Schritt, der einen numerischen Readout aus einer Konstruktion
heraus\-trägt, in einen mit einer zweiten Konstruktion geteilten Container hinein,
und dann am Ergebnis eine Ordnungsrelation abliest. Die Pflichten verlangen im
Umriss, dass beide Seiten explizite numerische Readouts tragen, dass der geteilte
Container mit typisierten Einbettungen kommt, dass die Identifikation zwischen
den Seiten eine erlaubte Identifikation und nicht eine Namensgleichheit ist, dass
das Indeterminacy-Budget in absorbierte und nichtabsorbierte Anteile zerlegt ist,
dass eine numerische Auswertung im Container existiert, und --- als letzte und
entscheidende --- dass die Ordnungsrelation aus der erlaubten Identifikation
\emph{abgeleitet} und nicht neben ihr behauptet wird. Operationalisiert sind die
Pflichten als Ledger aus neun Zeilen, jede eine Frage mit protokollierter
Erwartung und protokolliertem Risiko, gegen einen Primärtext zu füllen. Um das
Ledger herum liegen fünf Disziplinen, die seine Befunde nachprüfbar machen: ein
Verdikt-Vokabular, das Form von Wahrheit trennt, eine Typologie, die Lücken
klassifiziert statt sie zu graduieren, ein Import-Ledger für jeden Rückgriff auf
Standardliteratur, eine Regel für die Berichterstattung über vergeblich Gesuchtes
und eine Regel für die Offenlegung von Rekonstruktionen bei unvollkommener
Textbasis. Abschnitt~\ref{sec:instrument} führt all das aus.

\subsection{Quellen, Kürzel und Ankerkonvention}
\label{sec:sources}

Sieben Dokumente werden verwendet, in Gestalt von Textextrakten ihrer PDFs: sechs
tragen die vier Anwendungen, das siebte ist der in Abschnitt~\ref{sec:convergence}
geprüfte Zwischenbericht. Die Kürzel der linken Spalte gelten durchgehend; jede
Quellenangabe dieser Arbeit trägt eines, weil die Extrakte unabhängig
durchnummeriert sind und eine bloße Zeilennummer mehrdeutig wäre (siehe Tabelle~\ref{tab:primary_sources}).

\begin{table}[htbp]
\centering
\small
\renewcommand{\arraystretch}{1.12}
\begin{tabular}{@{}lll@{}}
\toprule
Kürzel & Quelle & Rolle in dieser Arbeit \\
\midrule
IUT-III & \cite{Mochizuki2021c} & Fall A: das Korollar und sein Beweis \\
IUT-IV & \cite{Mochizuki2021d} & Fall B: der deklarierte Quantifizierungsort \\
ATS-III & \cite{JoshiATSIII} & Fall C: Konstruktion und Volumentheorie \\
ATS-IV & \cite{JoshiATSIV} & Fall C: der $abc$-Schluss und seine Brücke \\
Zhou-I & \cite{ZhouI} & Fall D: erste nachgelagerte Anwendung \\
Zhou-II & \cite{ZhouII} & Fall D: zweite nachgelagerte Anwendung \\
LANA & \cite{LANA2026} & Abschnitt~\ref{sec:convergence}: unabhängige Formalisierung \\
\bottomrule
\end{tabular}
\caption{Primärquellen-Korpus und methodische Rollen in den vier Fall-Auditierungen.}
\label{tab:primary_sources}
\end{table}

Zitate stehen wörtlich so in den Extrakten; mathematische Symbole sind gesetzt,
Auslassungen sind markiert, und jede Rekonstruktion eines bei der Extraktion
verlorenen Symbols wird an Ort und Stelle offengelegt. Hervorhebungen innerhalb
von Zitaten stammen von uns, sofern nicht anders vermerkt.

\paragraph{Zitatsprache.} Die Zitate bleiben im englischen Originalwortlaut. Das
ist kein Versäumnis der Übersetzung, sondern eine Konsequenz des Gegenstands: Der
Wortlaut \emph{ist} hier die Evidenz. Ob ein Text \glqq follows formally\grqq{} schreibt oder
einen Satz beweist, ist genau die Frage; eine Übersetzung dieser Stellen würde
den Befund durch seine Wiedergabe ersetzen. Die Einordnung jedes Zitats steht
deutsch daneben.

\subsection{Aufbau}

Abschnitt~\ref{sec:instrument} führt das Instrument vollständig aus --- alle sieben
Pflichten und alle neun Zeilen --- und Abschnitt~\ref{sec:recipe} gibt das
Anwendungsrezept an, damit es gegen ein anderes Transportargument als die hier
geprüften vier ausgeführt werden kann. Die Abschnitte~\ref{sec:caseA}--\ref{sec:caseD} berichten die vier
Anwendungen. Abschnitt~\ref{sec:convergence} sagt, was sie gemeinsam zeigen --- und
ebenso ausführlich, was nicht. Abschnitt~\ref{sec:limitations} listet die
Einschränkungen, Abschnitt~\ref{sec:conclusion} schließt ab.

\paragraph{Ort in der Serie.} Dies ist Teil~3 der \emph{abc}-Gap-Serie. Teil~1 ist
die veröffentlichte forensische Untersuchung derselben Passage~\cite{GeigerGap2026},
auf der Abschnitt~\ref{sec:caseA} aufbaut und von der er dort abgegrenzt wird. Das
hier ausgeführte Zertifikat ist eine Instanz einer breiteren Tor-Methodik --- die
Instanz für numerische Membership ---, und diese breitere Methodik ist inzwischen
als Teil~2 der Serie~\cite{GeigerBCM2026} erschienen. Version 0.1 dieser Arbeit
setzte nichts daraus voraus, weil es sie noch nicht gab. Diese Version übernimmt
eine Sache daraus und sagt, welche: die gehärtete Lesart von Pflicht~6
(Abschnitt~\ref{sec:obligations}), derentwegen zwei Verdikte der ersten Version
hier umcodiert werden --- siehe \glqq Änderungen in Version 0.2\grqq{} oben.
Was diese Arbeit beiträgt, ist die Instanz und die vier unten berichteten
Anwendungen, und diese füllen das Ledger nur zum Teil: Fall~A berichtet vier
seiner neun Zeilen, wie die Abschnitte~\ref{sec:caseA} und~\ref{sec:convergence}
festhalten.

% =============================================================================
\section{Das Brücken-Zertifikat}
\label{sec:instrument}
% =============================================================================

\subsection{Die Frage}

Das Instrument betrifft Schritte folgender Gestalt. Zwei Konstruktionen erzeugen
zwei Objekte. Das erste trägt einen numerischen Readout $\nu_q$ --- in den unten
behandelten Fällen ein normiertes logarithmisches Volumen. Das zweite bestimmt
eine Zielmenge, in den unten behandelten Fällen eine abgeschlossene reelle
Halbgerade $H_\Theta$. Die beiden Objekte werden dann in einem gemeinsamen
Container vorgeführt, verbunden durch Daten, die wechselweise als Verknüpfung,
Identifikation, Transport oder Korrespondenz beschrieben werden, und eine Aussage
der Form
\[
  \nu_q(\text{erstes Objekt}) \in H_\Theta
\]
wird getroffen. Aus dieser Membership folgt unmittelbar eine Ungleichung, und
diese Ungleichung ist der Zweck der ganzen Konstruktion.

Die Frage des Zertifikats lautet:

\begin{quote}
Wann darf ein qualitatives Paket --- zwei Objekte, gemeinsame Labels, ein
algorithmischer Transport, ein gemeinsamer Container --- als numerische
Membership-Aussage in einer konkreten reellen Halbgeraden gelesen werden?
\end{quote}

Das ist bewusst enger als \glqq ist das Argument korrekt?\grqq{} und bewusst weiter als \glqq ist
diese Formel richtig?\grqq. Es ist die Frage, um die der Streit tatsächlich geht.

\subsection{Die sieben Pflichten}
\label{sec:obligations}

In der vorgeschlagenen Spezifikation enthält ein zulässiges Zertifikat $\Brxi$
mindestens Folgendes. Dies sind Pflichten des gewählten Formvertrags, keine
bewiesenen notwendigen Bedingungen für jeden möglichen mathematischen Beweis.

\begin{obligation}[$q$-seitiges Objekt mit Readout]
Ein Objekt $\mathrm{Obj}_q$ samt konkretem logvolumetrischem Readout $\nu_q$.
\end{obligation}

\begin{obligation}[$\Theta$-seitiges Objekt mit Ziel]
Ein Objekt $\mathrm{Obj}_\Theta$ samt Zielhalbgerade
$H_\Theta = \mathbb{R}_{\le -|\log(\Theta)|}$.
\end{obligation}

\begin{obligation}[Gemeinsamer Container mit typisierten Einbettungen]
Ein Container $C$ samt explizit typisierten Einbettungen $i_q$, $i_\Theta$.
\end{obligation}

\begin{obligation}[Linkage-Daten]
Label- und Linkage-Daten $L$, die nicht bloß angeben, dass beide Seiten dieselben
Namen tragen, sondern dass eine bestimmte Identifikation \emph{erlaubt} ist.
\end{obligation}

\begin{obligation}[Getrenntes Indeterminacy-Budget]
Ein Indeterminacy-Budget $I$, getrennt nach absorbierbarer Renormierung,
Blurring und nichtabsorbierter Restlast.
\end{obligation}

\begin{obligation}[Kompatible Auswertung im Container]
\label{obl:six}
Eine numerische Auswertung $\nu_C$ im Container \emph{zusammen mit} den
Kompatibilitätsidentitäten, die die Behauptung braucht --- insbesondere
$\nu_C \circ i_q = \nu_q$ auf dem ausgewiesenen Quellbereich. Diese
Kompatibilitätsklauseln sind zwingend: Die bloße Existenz einer numerischen
Abbildung auf $C$ löst diese Pflicht nie ein, und eine stipulierte Invarianz gilt
nicht als Beweis einer Kompatibilität.
\end{obligation}

\begin{obligation}[Ordnungskompatibilität]
\label{obl:seven}
Die Implikation selbst: Aus der erlaubten Linkage \emph{folgt}
$\nu_q(\mathrm{Obj}_q) \in H_\Theta$.
\end{obligation}

Die Last des Zertifikats sitzt fast vollständig in den beiden letzten Pflichten,
und ihre Unterscheidung ist das Nützlichste, was das Instrument liefert. Pflicht~6
fragt, ob die beiden Seiten durch eine ausgewiesene Kompatibilität auf eine Skala
gebracht sind --- nicht bloß, ob sie in einem Container untergebracht sind.
Pflicht~7 fragt, ob die Ordnungsrelation zwischen den beiden Auswertungen
abgeleitet ist. Die beiden können getrennt und aus verschiedenen Gründen
ausfallen, und sie auseinanderzuhalten ist das, was einen diffusen Verdacht in
zwei benennbare fehlende Aussagen verwandelt. Ohne die Pflichten~6 und~7 ist
Kohabitation in einem gemeinsamen Container qualitative Nähe, nicht die
numerische Aussage, die das Argument braucht.

\paragraph{Herkunft des Wortlauts der sechsten Pflicht.} Version 0.1 dieser
Arbeit formulierte Pflicht~6 disjunktiv: eine numerische Auswertung im Container
\emph{oder} eine beweisbare Kompatibilität. Die konjunktive Fassung oben ist dem
Methodenpapier der Serie~\cite{GeigerBCM2026} entnommen, in dem die entsprechende
Vertragskomponente nach dem Erscheinen dieser Arbeit gehärtet wurde. Die Änderung
ist nicht kosmetisch --- sie ist der Grund, weshalb hier zwei Verdikte aus
Version 0.1 umcodiert werden (\glqq Änderungen in Version 0.2\grqq) --- und sie löst eine
Spannung auf, die diese Arbeit von Anfang an trug: Das Arbeitslemma weiter unten
verlangte bereits einen \glqq labelkompatiblen, skalenerhaltenden\grqq{} Auswertungsvertrag,
also mehr, als die disjunktive Fassung von Pflicht~6 forderte.

\paragraph{Wann eine Pflicht als eingelöst gilt.} Weil die Pflichten von
verschiedener Art sind, ist auch die Evidenz verschieden, die sie einlöst; das
vorab zu sagen, hält die Codierung symmetrisch statt bequem. Die Pflichten~1--5
können durch eine ausdrückliche Festlegung eingelöst werden: eine Definition, eine
deklarierte Eigenschaft, eine fixierte Konvention. Pflicht~6 verlangt mehr als ein
gemeinsames Symbol und mehr als einen gemeinsamen Container: eine prüfbare
Zuordnung beider Seiten zu einer Auswertung \emph{und} die
Kompatibilitätsidentitäten, die diese Auswertung an die Readouts binden, die sie
darstellen soll. Dass beide Seiten reelle Zahlen sind, gehört nicht dazu, und
eine stipulierte Invarianz ebenso wenig. Pflicht~7 wird nur dort
als eingelöst protokolliert, wo der Text eine getrennt identifizierbare Herleitung,
ein benanntes Lemma oder ein ausdrücklich delegiertes Resultat enthält, das die
Ordnungsrelation aus der erlaubten Identifikation gewinnt.

Die negative Form dieser letzten Regel bindet diese Arbeit am stärksten, und sie
ist bewusst schwach: Ein solches Element unter dem berichteten Quellenumfang und
den berichteten Suchbegriffen nicht zu finden, ist ein \emph{dokumentarischer
Nichtbefund}. Es ist kein Nachweis, dass keine Herleitung existiert, und kein
Urteil darüber, ob die deklarierten Eigenschaften ausreichten, wenn die Herleitung
ausgeschrieben wäre. Diese zweite Frage ist mathematisch, und
Abschnitt~\ref{sec:cannot} hält fest, dass das Instrument sie nicht entscheidet. Tabelle~\ref{tab:obligations_taxonomy} bietet eine synoptische Taxonomie der sieben formalen Pflichten sowie der getriggerten Erweiterung, einschließlich ihrer mathematischen Objekte, operativen Anforderungen, zugeordneten Ledger-Zeilen und primären Fehlerrisiken.

\begin{table}[htbp]
\centering
\small
\setlength{\tabcolsep}{3.5pt}
\renewcommand{\arraystretch}{1.15}
\begin{tabular}{@{}lP{0.22\textwidth}P{0.25\textwidth}lP{0.22\textwidth}@{}}
\toprule
Pflicht & Mathematisches Objekt & Operative Anforderung & Zeile & Primäres Fehlerrisiko \\
\midrule
1 ($q$-Readout) & $(\mathrm{Obj}_q, \nu_q)$ & Expliziter log-volumetrischer Readout & R0, R2 & Skalenverlust / implizite Typisierung \\
2 ($\Theta$-Ziel) & $(\mathrm{Obj}_\Theta, H_\Theta)$ & Abgeschlossene Ziel-Halbgerade $H_\Theta \subset \mathbb{R}$ & R0, R3 & Nachträgliche Zielverschiebung \\
3 (Container) & $(C, i_q, i_\Theta)$ & Gemeinsamer Container mit Einbettungen & R1 & Fehlende Kohabitation \\
4 (Koppelung) & $L \subset \mathrm{Obj}_q \times \mathrm{Obj}_\Theta$ & Zulässige Identifikation (Protokoll) & R4 & Namenskoinzidenz ohne Erlaubnis \\
5 (Budget) & $I = (I_{\mathrm{abs}}, I_{\mathrm{res}})$ & Zerlegung in absorbierte/unabsorbierte Last & R5 & Verdeckter Zweitmodus \\
6 (Auswertung) & $(\nu_C, \nu_C \circ i_q = \nu_q)$ & Containermetrik mit zwingender Identität & R6a & Skalenloser Container / Identitätslücke \\
7 (Ordnung) & $\nu_q(\mathrm{Obj}_q) \in H_\Theta$ & Abgeleitete Ungleichungsimplikation & R6b & Behauptete Inklusion / Nicht-Ableitung \\
\midrule
Erw.~E8 & Modus-Zweig $M_2$ & Getriggerter sekundärer Auswertungskanal & E8 & Unausgewiesene Zweigwahl \\
\bottomrule
\end{tabular}
\caption{Taxonomie der sieben Brückenzertifikats-Pflichten, formalen Signaturen und Prüfkriterien.}
\label{tab:obligations_taxonomy}
\end{table}

\subsection{Das Ledger}
\label{sec:ledger}

Die Pflichten sind als neun Zeilen operationalisiert. Jede Zeile ist eine Frage
an den Wortlaut eines Primärtextes; jede trägt eine protokollierte Erwartung, die
\emph{vor} jeder Textprüfung festgelegt wurde, und ein protokolliertes Risiko ---
den Fehlermodus, dessentwegen die Zeile existiert. Die Erwartung vorab
festzuhalten ist keine Dekoration: Sie erlaubt einem späteren Leser zu sehen,
welche Befunde eine vorgängige Überzeugung bestätigt und welche ihr widersprochen
haben. Die Datierung gehört zur Behauptung und wird deshalb genannt: Die Pflichten,
die Zeilen und ihre Erwartungen wurden im Mai~2026 aufgeschrieben; die vier unten
berichteten Anwendungen wurden im August~2026 durchgeführt (Tabelle~\ref{tab:bc_ledger_rows}).

\begin{table}[htbp]
\centering
\small
\renewcommand{\arraystretch}{1.12}
\begin{tabular}{@{}lP{0.38\textwidth}P{0.17\textwidth}P{0.28\textwidth}@{}}
\toprule
Zeile & Test & Vorab-Erwartung & Risiko, das die Zeile abfängt \\
\midrule
R0 & Sind die beiden Linien getrennt typisiert? & ja & Scheingleichheit \\
R1 & Ist ein gemeinsamer Container vorhanden? & ja & Container ersetzt keine Skala \\
R2 & Bleibt der numerische Readout auf der $q$-Seite erhalten? & offen & Readout verliert seine Objektidentität \\
R3 & Steht die Zielhalbgerade auf der $\Theta$-Seite fest? & ja & Halbgerade wird nachträglich verschoben \\
R4 & Wird die $j^2$-Last adressiert? & offen & unverbuchte $O(\ell^2)$-Restlast \\
R5 & Ist das Indeterminacy-Budget getrennt? & teilweise & Blurring versteckt einen zweiten Modus \\
R6 & Wird Ordnungserhaltung bewiesen? & offen & qualitative Linkage reicht nicht \\
R7 & Ist Approximation schwächer als Membership? & bei geschlossener Halbgerade wahrscheinlich nein & Scheinausweg \\
R8 & Ist ein zweiter Modus sichtbar oder ausgeschlossen? & offen & der Kern der Frage \\
\bottomrule
\end{tabular}
\caption{BC-Ledger v0: Neun Prüfzeilen mit operativen Testfragen, Vorab-Erwartungen und erfassten Fehlerrisiken.}
\label{tab:bc_ledger_rows}
\end{table}

Drei Hinweise zum Lesen dieser Tabelle. Erstens heißt \glqq offen\grqq{} in der
Erwartungsspalte, dass keine Erwartung festgelegt wurde, nicht dass eine negative
Antwort erwartet wurde. Zweitens sind die Zeilen nicht gleichgewichtig. So wie sie
formuliert ist, benennt R6 nur die Ordnungsfrage, trägt aber tatsächlich die
Pflichten~6 und~7 zugleich --- eine Vermengung, die der Split in
Abschnitt~\ref{sec:splits} auflöst und die zu benennen sich lohnt: Solange eine
Zeile für die Ordnungsfrage existiert und keine für die Kompatibilitätsfrage,
wird die zweite im Fließtext statt am Ledger beurteilt. R6 ist die Zeile, an der
das Instrument informativ sein soll; R0, R1 und R3 sind Hygienezeilen, deren
Hauptfunktion darin besteht, einen negativen Befund bei R6 interpretierbar statt
generisch zu machen.

Drittens ist die letzte Zeile keine Kernzeile, und sie als solche zu lesen
verzerrt den Nenner. R8 fragt nach einem zweiten Auswertungsmodus, und das ist
ein Merkmal des Gegenstandsmodells, nicht des Vertrags: Sie bildet auf keine
Pflicht ab; sie ist nur dort aktiv, wo der Gegenstandsbereich der Behauptung
tatsächlich mehrere Modi, Zweige oder Normierungen anbietet; und wo dieser
Auslöser fehlt, ist sie außerhalb des Geltungsbereichs statt negativ. Das
Nachfolgeinstrument führt sie deshalb außerhalb des Pflichten-Crosswalks, als
getriggerte Erweiterung~E8, und berichtet einen Befund an diesem Label und nie an
einer Vertragskomponente~\cite{GeigerBCM2026}. An der unten protokollierten
Codierung ändert sich nichts; es ändert sich ihr Status --- und dass eine
Vollständigkeitsaussage die acht Kernzeilen R0--R7 meint.

\paragraph{Versionsetikett und Zeilen-Namespace.} Das hier ausgeführte Ledger ist
das im Mai~2026 fixierte Zertifikats-Ledger. Seine Zeilen sind präfixlos, was
harmlos war, solange es das einzige Ledger der Serie war, und nicht mehr harmlos
ist: Das Nachfolgeinstrument~\cite{GeigerBCM2026} versieht seine Zeilen mit einem
Versionspräfix (v1-T für das Transport-Ledger, v1-F für das vollständige), und
gleiche Ziffernsuffixe über Instrumente hinweg bedeuten \emph{nicht} gleiche
Semantik --- das R8 dieser Arbeit ist die Zweitmodus-Zeile, während die Zeile mit
derselben Nummer im Transportinstrument eine andere Frage stellt. Jedes präfixlose
Zeilenlabel im Folgenden ist ein Label des Ledgers dieser Arbeit, und die
Entsprechung zu den Nachfolgezeilen ist ein Crosswalk, keine Identität. Wo ein
Versionsetikett gebraucht wird, weist diese Version dem Ledger den Namen
\textbf{BC-Ledger~v0} zu; das Etikett wird hier eingeführt und stammt nicht aus
dem Serienregister, das für dieses Ledger bislang keines führt.

\subsection{Verfeinerungen aus dem Gebrauch}
\label{sec:splits}

Drei Zeilen erwiesen sich als Vermischungen verschiedener Fragen, und das
Instrument wurde entsprechend geschärft. Die ersten beiden Verfeinerungen sind
\emph{Ergebnisse der Anwendungen} und waren nicht Bestandteil des fixierten
Zertifikats; die dritte folgt aus der oben ausgeführten gehärteten Lesart von
Pflicht~6. Keine der drei ist noch bloße Empfehlung: Im Nachfolgeinstrument sind
alle drei normativ~\cite{GeigerBCM2026}. Die Fallberichte der
Abschnitte~\ref{sec:caseA}--\ref{sec:caseD} bleiben überall dort bei den
\emph{ursprünglichen} Zeilen, wo der Split am Befund nichts ändert, und berichten in der Regel
am Split, wo er etwas ändert --- Fall~B an R4a und R4b, Fall~A an R6a und R6b.
Auch die Befunde von Fall~C trennen sich an den beiden Hälften; sein Bericht
behält die zusammengelegte Zeile der Version 0.1 bei, die Trennung wird im
Fließtext von Abschnitt~\ref{sec:caseC} geführt.

\paragraph{R2a / R2b.} Zeile R2 mischt zwei Transporte. \emph{R2a} fragt, ob der
Readout den Transport \emph{innerhalb} eines Containers übersteht --- entlang des
Links, den die Konstruktion selbst bereitstellt. \emph{R2b} fragt, ob er den
Übergang \emph{zwischen} Containern übersteht. Nur R2b ist die eigentliche
Brückenfrage. Die Unterscheidung wurde durch die vergleichende
Diskriminierungsanwendung in
Abschnitt~\ref{sec:caseC} erzwungen, wo die beiden Transporte sichtbar
verschiedener Art sind, und durch Abschnitt~\ref{sec:caseA} bestätigt, wo R2b mit
R6 zusammenfällt, weil die Identifikation der Container \emph{der} Brückenschritt
ist.

\paragraph{R4a / R4b.} Zeile R4 mischt den Ort, an dem eine Last beziffert wird,
mit dem Ort, an dem sie gegen die Brücke verbucht wird. \emph{R4a} fragt, ob die
Last in der Auswertung der $\Theta$-Seite überhaupt beziffert wird. \emph{R4b}
fragt, ob sie an der Brückenstelle selbst verbucht wird. Die Unterscheidung wurde
durch Abschnitt~\ref{sec:caseB} erzwungen: R4a ist mit ja zu beantworten und R4b
mit nein, und ein einziges Wort für beides verdeckt zwei verschiedene Befunde.

\paragraph{R6a / R6b.} Zeile R6 mischt die beiden tragenden Pflichten unter einer
Frage, die nur eine von beiden benennt. \emph{R6a} fragt nach Pflicht~6: Ist die
gemeinsame Auswertung von den Kompatibilitätsidentitäten begleitet, die sie an die
beiden Readouts binden? \emph{R6b} fragt nach Pflicht~7: Ist die Ordnungsrelation
aus der erlaubten Identifikation abgeleitet? Dieser Split ist nicht durch die
Anwendungen entstanden, sondern durch die gehärtete Lesart von Pflicht~6
(Abschnitt~\ref{sec:obligations}). Sobald eine gemeinsame Auswertung die sechste
Pflicht nicht mehr für sich allein einlöst, kann eine einzige Zeile nicht beide
Verdikte tragen --- Abschnitt~\ref{sec:caseC} gibt an den beiden Hälften
verschiedene Antworten zurück, und Abschnitt~\ref{sec:caseA} gibt dieselbe
Antwort aus zwei verschiedenen Gründen; in beiden Konstellationen hätte
eine einzige zusammengelegte Zeile die Struktur verdeckt.

Eine weitere Verfeinerung wird als offener Punkt vermerkt statt übernommen:
Pflicht~3 (Container-Typisierung) geht bislang nur über R1 ins Ledger ein, und R1
fragt, ob ein Container \emph{vorhanden} ist. Sowohl
Abschnitt~\ref{sec:caseA} als auch Abschnitt~\ref{sec:caseC} legen nahe, dass die
interessante Frage lautet, ob die Typisierung \emph{bis zum Schluss durchgehalten}
wird --- eine andere Frage, die eine eigene Zeile verdient.

\subsection{Verdikte und Konfidenz}

Eine Zeile erhält eines von drei Verdikten, und das Vokabular ist festgelegt, weil
lockere Verdiktsprache der Weg ist, auf dem Formbefunde zu Wahrheitsbehauptungen
werden (siehe Tabelle~\ref{tab:verdict_taxonomy}).

\begin{table}[htbp]
\centering
\small
\renewcommand{\arraystretch}{1.12}
\begin{tabular}{@{}lP{0.70\textwidth}@{}}
\toprule
Verdikt & Bedeutung \\
\midrule
\textsc{liefert} & Die Pflicht wird vom geprüften Wortlaut eingelöst. \\
\textsc{teilweise} & Ein Teil der Pflicht wird eingelöst und ein benannter Teil
  nicht. Die Aufteilung ist zu nennen, nicht zu mitteln. \\
\textsc{liefert nicht} & Die Pflicht wird vom geprüften Wortlaut nicht eingelöst.
  Das ist \emph{keine} Behauptung, die Aussage sei falsch, und keine, sie könne
  nicht anderswo eingelöst werden. \\
\bottomrule
\end{tabular}
\caption{Verdikt-Vokabular für Zertifikats-Beurteilungen.}
\label{tab:verdict_taxonomy}
\end{table}

Jedes Verdikt trägt eine Konfidenz, und die Konfidenz betrifft den
\emph{Textbefund}, nicht die Mathematik: Hohe Konfidenz heißt, dass die Zitate
eindeutig sind und dass kein Verdikt an einer zweifelhaften Lesart hängt. Ein
Verdikt kann zugleich hochkonfident und in seiner Reichweite eng sein; wo das
vorkommt, steht die Reichweitenbeschränkung neben dem Verdikt statt in ihm.

Diese Stufen sind autorseitige Einschätzungen der Textlage --- Eindeutigkeit des
Zitats, Genauigkeit des Ankers, Abdeckung der Suche und getragene
Rekonstruktionslast. Sie sind keine kalibrierten Wahrscheinlichkeiten, und es gab
keine unabhängige Zweitcodierung, also wird auch keine Übereinstimmung zwischen
Codierenden berichtet. Sie sind als Aussage darüber zu lesen, wie fest der
Wortlaut den Eintrag trägt, nicht als gemessene Zuverlässigkeit.

\subsection{Lücken-Typologie}
\label{sec:typology}

Lücken werden typisiert, nicht graduiert. Graduierung (\glqq schwerwiegend\grqq,
\glqq geringfügig\grqq) schmuggelt eine Bewertung ein, zu der das Instrument nicht befugt
ist; Typisierung sagt, welche Art von Sache fehlt, was zugleich nützlicher und
besser verteidigbar ist (Tabelle~\ref{tab:gap_taxonomy}).

\begin{table}[htbp]
\centering
\small
\renewcommand{\arraystretch}{1.12}
\begin{tabular}{@{}lP{0.76\textwidth}@{}}
\toprule
Typ & Beschreibung \\
\midrule
G & Ein fehlender oder nicht nachvollziehbarer Beweisschritt. \\
Z & Ein Zitat trägt die behauptete Aussage nicht --- ein \emph{Import}-Fehler. \\
D & Definitionsdrift: Ein Begriff wird anders benutzt, als er definiert wurde. \\
N & Kein Fehler, aber eine beweisrelevante Unschärfe auf Konventions- oder
  Notationsebene. \\
\bottomrule
\end{tabular}
\caption{Typologie identifizierter Lücken und Textdefizite.}
\label{tab:gap_taxonomy}
\end{table}

Die Typen kombinieren. Ein Schritt, dessen Reparatur sowohl ein fehlendes
Argument als auch die Auflösung einer definitorischen Mehrdeutigkeit verlangt,
wird als G+D erfasst, und beide Teile sind zu benennen.

\subsection{Import-Ledger-Disziplin}
\label{sec:importledger}

Jeder Rückgriff auf Standardliteratur wird als \emph{Import} registriert, mit
exakter Fundstelle in der publizierten Quelle. Die zugehörige Regel ist bewusst
asymmetrisch:

\begin{quote}
Importe aus peer-reviewter Literatur gelten als vertrauenswürdig. Geprüft wird
ihre \emph{Passung} --- ob die zitierte Aussage die Aussage ist, die das Argument
braucht --- nicht ihr Beweis.
\end{quote}

Das ist der operative Gehalt des Z-Typs. Die Regel hält die Prüfung endlich: Ohne
sie entartet jede Prüfung zu einer Prüfung der gesamten arithmetischen Geometrie.
Sie hält die Prüfung auch in der Gegenrichtung ehrlich, denn ein Import, dessen
Passung nicht geprüft wurde, muss als ungeprüft vermerkt und darf nicht
übergangen werden.

\subsection{Belegregeln}
\label{sec:evidence}

Vier Regeln bestimmen, was in eine Ledger-Zeile geschrieben werden darf.

\paragraph{Kein Autoritätsargument.} Weder die Behauptung des geprüften Autors,
ein Schritt genüge, noch die Behauptung eines Kritikers, er genüge nicht, ist
Beleg. Das gilt in beide Richtungen und ist keine Höflichkeit: Die
Selbsteinschätzung eines Autors ist Prüfgegenstand, und das Verdikt eines
Kritikers ist für das Ledger nur die Definition dessen, wie eine Negativkontrolle
aussähe.

\paragraph{Wörtliches Zitat mit Anker.} Jeder Befund wird von einem zitierten
Satz mit Zeilenanker in einen benannten Extrakt getragen. Ein Befund, der nicht
zitiert werden kann, ist kein Befund.

\paragraph{\glqq Gesucht, nicht gefunden\grqq.} Die Behauptung, etwas fehle in einem Text,
wird zusammen mit den verwendeten Suchbegriffen und dem durchsuchten Bereich
berichtet. Diese Regel macht negative Befunde falsifizierbar: Wer die Suche für
zu eng hält, kann genau sehen, wie eng sie war, und sie wiederholen. Jeder
Fallbericht, der auf einer Abwesenheitsbehauptung ruht, trägt einen solchen Block:
Die Fälle~A, B und~D unten tun das. Fall~C behauptet keine Abwesenheit --- seine
Befunde ruhen auf zitierten Passagen, die vorhanden sind --- und trägt deshalb
keinen.

\paragraph{Rekonstruktionsdisziplin.}
\label{sec:reconstruction}
Aus PDF extrahierter Text verliert Relationssymbole, griechische Buchstaben,
Exponenten und Indizes. Wo ein lasttragender Satz betroffen ist, sind drei Dinge
verlangt: Die Rekonstruktion muss ausdrücklich genannt, ihre Herleitung aus
anderen Passagen desselben Dokuments angegeben und eine Gegenprobe an
unabhängiger Stelle vorgeführt werden. Wo die naive Lesart einer beschädigten
Zeile den Sinn umkehren würde, ist diese Zeile eigens zu markieren. Die Regel ist
Bestandteil des Instruments, keine Entschuldigung dafür --- und sie ist hier
relevant, weil in einer der vier Anwendungen ein Verdikt tatsächlich an einer
solchen Rekonstruktion hängt (Abschnitt~\ref{sec:reconstruction-A}).

\subsection{Das Arbeitslemma}

Das ganze Instrument ist in einer Aussage verdichtet, die eine Aufgabenbeschreibung
und kein Satz ist.

\begin{workinglemma}
Für die forensische Lesart eines Schritts der in
Abschnitt~\ref{sec:instrument} beschriebenen Gestalt genügt es nicht, eine
qualitative Relation zwischen den beiden Objekten in einem gemeinsamen Container
zu zeigen. Erforderlich ist ein expliziter, labelverträglicher, skalenbewahrender
und ordnungskompatibler numerischer Auswertungsvertrag $\Brxi$. Jeder Kandidat
für einen zweiten Modus $T_2$ muss entweder in diesem Vertrag sichtbar sein oder
als unzulässige versteckte Hilfsskala verbucht werden.
\end{workinglemma}

Durchgehend bezeichnet $T_2$ \emph{den zweiten Auswertungsmodus} --- eine zweite
Skala, einen Shift oder eine zweite Verflechtung neben derjenigen, die das
Argument ausweist --- und nichts sonst. Zeile R8 ist die Ledger-Zeile, die nach
ihm fragt.

Der Status dieser Aussage lautet: Arbeitslemma, nicht bewiesen. Sie ist als
präziser Prüfauftrag für Primärtextstellen und für Positiv- und Negativkontrollen
zu behandeln --- was der Rest dieser Arbeit mit ihr tut.

\subsection{Was das Instrument nicht entscheiden kann}
\label{sec:cannot}

Positiv formuliert, damit es nicht als Bescheidenheit missverstanden wird.

\begin{enumerate}
\item \textbf{Es kann Stichhaltigkeit nicht bescheinigen.} Alle neun Zeilen können
  liefern und das Argument kann dennoch aus einem Grund scheitern, nach dem keine
  Zeile fragt.
\item \textbf{Es kann mathematische Sachfragen nicht entscheiden.} Ob ein Paar
  deklarierter Eigenschaften als Vertragslemma genügte, wenn es bewiesen wäre,
  oder ob eine identische globale Normierung zweier isomorpher Kopien eines
  Zahlkörpers als beweisbare Kompatibilität durchgeht --- das sind mathematische
  Fragen. Das Ledger hält nur fest, dass der Text ein Eigenschaftspaar oder eine
  Konvention liefert, wo er ein Lemma bräuchte.
\item \textbf{Es kann nicht über die geprüften Dokumente hinaussehen.} Eine Zeile
  mit \glqq liefert nicht\grqq{} ist eine Aussage über einen Wortlaut. Die Pflicht kann in
  ungeprüftem Material eingelöst sein; wo diese Möglichkeit konkret ist, benennen
  die Fallberichte sie.
\item \textbf{Es kann eine nicht eingelöste Pflicht nicht von einer uneinlösbaren
  unterscheiden.} Nichts am Instrument erlaubt den Schritt von \glqq hier nicht
  abgeleitet\grqq{} zu \glqq nicht ableitbar\grqq.
\end{enumerate}

\subsection{Anwendung auf ein neues Argument}
\label{sec:recipe}

Das Instrument ist nicht an die vier unten behandelten Fälle gebunden. Das
folgende Rezept ist ein aus den Anwendungen gewonnenes prospektives Protokoll.
Es wurde nicht in jedem berichteten Fall vollständig ausgeführt: Insbesondere
bleiben in Fall A Hygienezeilen offen, und die Hygiene-vor-R6-Regel entstand
erst nach den Anwendungen. Das Rezept wird für neue Transportargumente angeboten;
die historischen Teilanwendungen erfüllen dadurch nicht rückwirkend alle Schritte.

\begin{recipe}
\begin{enumerate}
\item \textbf{Den Schritt lokalisieren.} Den Satz identifizieren, an dem eine
  qualitative Relation zu einer numerischen Behauptung wird. Er steht in der
  Regel nicht dort, wo die Selbstzusammenfassung des Dokuments ihn vermutet; die
  Struktur am Text verifizieren und die Korrektur protokollieren, damit der
  nächste Leser sie nicht neu herleiten muss.
\item \textbf{Die beiden Seiten und den Container benennen.} Pflichten 1--3
  füllen. Die Typisierung der Einbettungen protokollieren --- und protokollieren,
  ob sie irgendwo vor dem Schluss fallengelassen wird.
\item \textbf{Die Hygienezeilen} R0, R1, R3 \textbf{füllen.} Fällt eine von ihnen
  aus, kann ein späterer Befund bei R6 ein Artefakt dieses Ausfalls statt ein
  eigenständiges Ergebnis sein.
\item \textbf{Pflicht~6 von Pflicht~7 trennen.} Zuerst fragen, ob der Container
  eine numerische Auswertung trägt \emph{und} die Kompatibilitätsidentitäten, die
  sie an die beiden Readouts binden --- insbesondere $\nu_C \circ i_q = \nu_q$
  auf dem ausgewiesenen Quellbereich. Eine gemeinsame Auswertung allein löst
  Pflicht~6 nie ein, und eine stipulierte Invarianz ist kein Beweis einer
  Kompatibilität. Erst dann fragen, ob die Ordnungsrelation aus der erlaubten
  Identifikation abgeleitet ist. Beide Fragen zu vermengen ist die häufigste Art,
  einen Befund zu erzeugen, der zugleich unfair und uninformativ ist; die erste
  auf die bloße Existenz einer gemeinsamen Skala zusammenzuziehen ist die
  häufigste Art, einen zu großzügigen zu erzeugen.
\item \textbf{Nach dem Vertragslemma ausdrücklich suchen}, mit protokollierten
  Suchbegriffen, über das ganze Dokument und getrennt über den Beweisbereich. Das
  Ergebnis in beiden Fällen festhalten.
\item \textbf{Jede gefundene Lücke typisieren} (Abschnitt~\ref{sec:typology}) und
  jeden Rückgriff auf externe Literatur ins Import-Ledger eintragen
  (Abschnitt~\ref{sec:importledger}).
\item \textbf{Den Block \glqq gesucht, nicht gefunden\grqq{} vor dem Verdikt schreiben},
  nicht danach. Seine Funktion ist, das Verdikt zu beschränken, und das kann er
  nicht, wenn er passend zu einem Verdikt verfasst wurde.
\item \textbf{Die Reichweite getrennt vom Verdikt angeben.} \glqq Liefert nicht, hohe
  Konfidenz, beschränkt auf Dokument $X$\grqq{} ist ein vollständiger Befund; \glqq liefert
  nicht\grqq{} allein ist keiner.
\end{enumerate}
\end{recipe}

% =============================================================================
\section{Fall A: das Korollar und sein Beweis}
\label{sec:caseA}
% =============================================================================

Die erste Anwendung gilt der Passage, für die das Instrument gebaut wurde: dem
Beweis von Corollary 3.12 im dritten Band der interuniversellen
Teichmüller-Theorie~\cite{Mochizuki2021c}. Gefüllt wurden die Zeilen R2, R4, R6
und R8 --- R6 an beiden Hälften des in Abschnitt~\ref{sec:splits} eingeführten
Splits, und R8 als die getriggerte Erweiterungszeile, die sie ist, nicht als
Kernzeile; R0, R1, R3, R5 und R7 waren nicht Teil dieser Anwendung und bleiben
offen (Abschnitt~\ref{sec:convergence}).

\paragraph{Verhältnis zu Teil~1 der Serie.} Diese Passage ist bereits gedruckt
lokalisiert. Eine frühere, reparaturorientierte Untersuchung desselben
Programms~\cite{GeigerGap2026} verortet die Lücke am Übergang zwischen denselben
Substeps und berichtet denselben Wortlaut --- dass die verlangte Ungleichung
\glqq formal folge\grqq. Die vorliegende Anwendung wiederholt diese Diagnose nicht und
verstärkt sie nicht, und sie beansprucht ihr gegenüber keine Priorität. Sie leistet
Engeres: Sie codiert die dokumentarische Frage als eine Pflicht eines ausgewiesenen
Zertifikats um, damit derselbe Ort mit drei weiteren Texten unter einem
unveränderten Instrument verglichen werden kann. Dass zwei Wege verschiedener Art
an derselben Adresse ankommen, wird hier als Befund über die Adresse festgehalten,
nicht als zweite Detektion.

Dies ist somit eine \emph{Teil}anwendung des Ledgers, und sie wird als solche
berichtet. Das Rezept aus Abschnitt~\ref{sec:recipe}, das verlangt, die
Hygienezeilen R0, R1 und~R3 zu füllen, bevor R6 gewertet wird, entstand später,
aus der Erfahrung dieser Anwendungen --- ebenso wie die Zeilenaufspaltungen aus
Abschnitt~\ref{sec:splits}, und aus demselben Grund werden die Fallberichte nicht
nachträglich daran angepasst. Die Folge wird benannt statt aufgesogen: Was folgt,
ist ein Befund über vier Zeilen. Die zu jeder Zeile protokollierte Konfidenz ist
Konfidenz in diese Zeile; ein Gesamturteil des Zertifikats über diese Passage wird
hier nicht beansprucht.

\subsection{Die Passage lokalisieren}

Die Substeps, in denen der Übergang stattfindet, stehen \emph{im Beweis} des
Korollars, nicht in der Bemerkung, die ihn rekapituliert. Die Struktur, hier
festgehalten, damit sie nicht neu hergeleitet werden muss (Tabelle~\ref{tab:case_a_text_blocks}):

\begin{table}[htbp]
\centering
\small
\renewcommand{\arraystretch}{1.12}
\begin{tabular}{@{}lP{0.30\textwidth}@{}}
\toprule
Block & Auszug Zeilen \\
\midrule
Proposition 3.9 (Log-Volumen für Pakete und Prozessionen) & 6365\,ff. \\
Theorem 3.11 (multiradiale Algorithmen), Beweis & 8666\,ff., Beweis 9032 \\
Remark 3.11.1 mit den deklarierten Eigenschaften & 9036--9300 \\
Corollary 3.12, Aussage & 9851--9884 \\
Corollary 3.12, Beweis & 9886\,ff. \\
Step (x) (Indeterminacy-Bilanz) & 10237--10283 \\
\textbf{Substeps (xi-a)--(xi-h)} & \textbf{10285--10496} \\
Remark 3.12.2 einschließlich Toy-Modell & 10573\,ff., Toy-Modell 10733--10813 \\
\bottomrule
\end{tabular}
\caption{Strukturblöcke und Zeilenanker des IUT-III Korollar-3.12-Beweisauszugs.}
\label{tab:case_a_text_blocks}
\end{table}

Vier in Remark 3.11.1 deklarierte Eigenschaften kehren durchgehend wieder und
werden wie in der Quelle abgekürzt: input prime-strip link (IPL),
strip-holomorphe Ausdrückbarkeitsbedingung (SHE), algorithmic parallel transport
(APT) und hidden internal structures (HIS).

\subsection{Offenlegung der Rekonstruktion}
\label{sec:reconstruction-A}

Der Extrakt dieses Dokuments verliert Relationssymbole, Exponenten, Indizes und
griechische Buchstaben. Anders als in den drei übrigen Anwendungen hängt hier
\emph{ein} Verdikt an einer Rekonstruktion; die Herleitung wird deshalb
vollständig angegeben, gemäß Abschnitt~\ref{sec:reconstruction}.

\begin{enumerate}
\item Die Zielmenge wird bei Z.~10385 definiert; der Extrakt liest
  \texttt{R-|log()| d=ef \{ R | -|log()|\} R}, zu lesen als
  $\mathbb{R}_{\le -|\log(\Theta)|} := \{\lambda \in \mathbb{R} \mid
  \lambda \le -|\log(\Theta)|\} \subset \mathbb{R}$.
\item Die Aussage des Korollars selbst legt die Richtung der Relation fest, bei
  Z.~9882/9884: \glqq \dots\ i.e., $C_\Theta \ge -1$ for any real number
  $C_\Theta \in \mathbb{R}$ such that
  $-|\log(\Theta)| \le C_\Theta \cdot |\log(q)|$\grqq, also
  $-|\log(\Theta)| \ge -|\log(q)|$.
\item Beides zusammen macht Z.~10446/10448 eindeutig:
  $-|\log(q)| \in \mathbb{R}_{\le -|\log(\Theta)|}$, \glqq hence also the inequality\grqq
  $-|\log(q)| \le -|\log(\Theta)| \in \mathbb{R}$.
\end{enumerate}

Die strittige Membership steht damit als \emph{ausdrückliche Aussage} bei
Z.~10446 --- mit den oben wiederhergestellten Relationssymbolen --- und ist nicht
interpretativ von uns erschlossen.

\paragraph{Prüfung am Seitenbild.} Weil ein Verdikt dieses Abschnitts an der
Wiederherstellung hängt, wurden die beiden betroffenen Zeilen nicht nur am Extrakt,
sondern an den veröffentlichten Seiten selbst geprüft. Beide Wiederherstellungen
sind bestätigt. Auf S.~184 trägt die Seite die Membership
$-|\log(q)| \in \mathbb{R}_{\le -|\log(\Theta)|}$ und darunter abgesetzt
$-|\log(q)| \le -|\log(\Theta)| \in \mathbb{R}$, gefolgt von \glqq then follows
formally\grqq; auf S.~174 trägt die Aussage des Korollars
$-|\log(\Theta)| \ge -|\log(q)|$ samt dem oben zitierten Satz. Die vom Extrakt
verlorenen Symbole stehen also in der Quelle, und zwar in der Form, die die
Rekonstruktion angenommen hat. Dabei wurde zugleich ein Detail unserer eigenen
Transkription korrigiert: Die Quelle schreibt durchgehend $C_\Theta$, nicht ein
bloßes $C$; die Zitate hier folgen ihr jetzt.

Eine weitere Zeile wird markiert, weil die naive Lesart ihren Sinn umkehrt. Bei
Z.~10750 liest der Extrakt \texttt{in the sense that R -h = -2h R}; gemeint ist
$\mathbb{R} \ni -h \neq -2h \in \mathbb{R}$, das \glqq $\neq$\grqq{} ist verloren. Der
umgebende Satz sichert das ab: Die beiden Zuordnungen seien \glqq mutually
incompatible and cannot be considered simultaneously\grqq
\anchor{IUT-III}{10748--10749}. Wer hier \glqq $=$\grqq{} einsetzt, schließt das Gegenteil
dessen, was die Passage sagt --- und das Seitenbild entscheidet es: S.~190 trägt
\glqq [in the sense that $\mathbb{R} \ni -h \neq -2h \in \mathbb{R}$]\grqq.

\subsection{R2 --- bleibt der \texorpdfstring{$q$}{q}-seitige Readout erhalten?}

\textsc{Liefert}, mit hoher Konfidenz --- und der Grund ist strukturell
interessant: Am entscheidenden Schritt findet gar kein Transport statt.

Die $q$-Seite wird in der Aussage des Korollars selbst als indeterminacy-frei
deklariert:

\begin{quote}\small
\glqq for the procession-normalized mono-analytic log-volume of the image of a
$q$-pilot object \dots\ in the multiradial representation of Theorem 3.11, (i),
\emph{which we do not regard as subject to the indeterminacies} (Ind1), (Ind2),
(Ind3) described in Theorem 3.11, (i), (ii).\grqq
\hfill\anchor{IUT-III}{9869--9873}
\end{quote}

Der unmittelbar vorangehende Absatz schreibt der $\Theta$-Seite genau das
Gegenteil zu \anchor{IUT-III}{9862--9864}. Die Asymmetrie ist also
\emph{deklariert}, nicht versehentlich. Der $q$-Pilot ist überdies der fixierte
Bezugspunkt der Konstruktion --- Eigenschaft (IPL) nennt ihn \glqq the
`coric'/`fixed' $q$-pilot $\mathcal{F}^{\Vdash\times}$-prime-strip\grqq
\anchor{IUT-III}{9117--9122} ---, und Eigenschaft (APT) hält fest, dass es einen
mengentheoretischen Transport gar nicht gibt:

\begin{quote}\small
\glqq This parallel transport mechanism \emph{does not consist of a simple instance of
transport of some set-theoretic region} \dots\ \emph{via some set-theoretic
function}. Rather, it consists of a construction algorithm that is simultaneously
valid/executable/well-defined with respect to the arithmetic holomorphic
structures in the domain and codomain of the $\Theta$-link\grqq
\hfill\anchor{IUT-III}{9212--9219}
\end{quote}

Der Readout ist konkret: $-|\log(q)|$ ist \glqq the log-volume of the region that
arises from the representation of the $q$-pilot object\grqq
\anchor{IUT-III}{10387--10390}, auf einer Skala, die durch die Aussage verankert
ist, \glqq global arithmetic degrees of objects of global realified Frobenioids may be
interpreted as log-volumes\grqq{} \anchor{IUT-III}{10391--10393}.

Pflicht~1 ist damit eingelöst, und Pflicht~4 ist auf der $q$-Seite eingelöst ---
aber \emph{stipulativ}: Die Pilot-Identität wird durch die Deklarationen
\glqq coric\grqq{}/\glqq fixed\grqq{} und \glqq do not regard as subject to\grqq{} geschützt, nicht durch eine
bewiesene Relation hergestellt.

\subsection{R4 --- wird die \texorpdfstring{$j^2$}{j\texttwosuperior}-Last adressiert?}

\textsc{Teilweise}, mit hoher Konfidenz. Der Buchungsort ist benannt; die Last
selbst wird dort nicht beziffert, und ihre Bezifferung ist ausdrücklich in den
vierten Band ausgelagert.

Die Last existiert und wird benannt: Die einschlägigen Monoide werden von den
Theta-Werten $\{q^{j^2}\}_{j=1,\dots,\ell^*}$ erzeugt
\anchor{IUT-III}{654--656}. Der Buchungsort ist die Prozessions-Normierung, und
verbucht werden dort \emph{Labels}:

\begin{quote}\small
\glqq the procession-normalized mono-analytic log-volumes [i.e., where the average is
taken over $j \in \mathbb{F}_\ell^*$] \dots\ are \emph{invariant with respect to
the indeterminacies} (Ind1), (Ind2), and have the effect of \emph{converting the
indeterminacy} (Ind3) \emph{into an inequality [from above]}.\grqq
\hfill\anchor{IUT-III}{10261--10268}
\end{quote}

wobei die Mittelung \glqq the $\mathbb{F}_\ell^*$-symmetry acting on the labels of the
theta values\grqq{} zugeschrieben wird \anchor{IUT-III}{10273--10276}. Das ist eine
Buchung von Symmetrie, nicht von Größenordnung.

Für das Korollar wird die Last dann ausdrücklich ausgeblendet, und die
Auslagerung wird offen deklariert. Eigenschaft (HIS) benennt genau diese
Anwendung:

\begin{quote}\small
\glqq if one takes the point of view --- \emph{as will be done in Corollary 3.12
below!} --- that one is only interested in considering the \emph{qualitative
logical aspects/consequences} \dots\ then \dots\ one may \dots\ \emph{forget that
this construction algorithm \dots\ has anything to do with theta functions \dots\
or theta values} $[\text{i.e., the } q^{j^2}]_{j=1,\dots,\ell^*}$!\grqq
\hfill\anchor{IUT-III}{9225--9244}
\end{quote}

und unmittelbar danach:

\begin{quote}\small
\glqq unlike Corollary 3.12 below, which only concerns qualitative logical
aspects/consequences \dots, \emph{the explicit computation to be performed in}
[IUT-IV], \S\,1, \emph{of the log-volumes that occur in the statement of Corollary
3.12 makes essential use of the way in which theta values occur} in the
construction algorithm of Theorem 3.11.\grqq{} \hfill\anchor{IUT-III}{9263--9267}
\end{quote}

Eine Suche über den gesamten Beweisbereich nach den Theta-Werten, der
Prozessions-Normierung und der Mittelung liefert Treffer \emph{ausschließlich} in
Step (x); in den Substeps selbst kommt die Last nicht mehr vor.

Ein Punkt der Fairness gehört ins Verdikt und nicht daneben. Der Fehlermodus,
dessentwegen diese Zeile existiert --- eine Last, die in einer Koordinate
verschwindet und unverbucht zurückkehrt ---, \emph{tritt hier nicht ein}. Was
eintritt, ist eine deklarierte Arbeitsteilung zwischen einem qualitativen und
einem quantitativen Band. Für das Zertifikat bleibt der Befund, dass die Last
innerhalb des geprüften Dokuments an der strittigen Stelle nicht verbucht ist; die
Unterscheidung zwischen \glqq anderswo verbucht, per Deklaration\grqq{} und \glqq verschleiert\grqq
wird zugunsten der Quelle festgehalten.

\subsection{R6 --- wird Ordnungserhaltung bewiesen?}

\textsc{Liefert nicht} an beiden Hälften der Zeile, mit hoher Konfidenz. Das ist
der Kernbefund der Anwendung, und er ist zweiteilig zu formulieren, weil die
beiden Pflichten aus verschiedenen Gründen ausfallen und die Zeile in ihrer
ursprünglichen Formulierung nur die zweite von beiden benennt.

\paragraph{R6a --- Pflicht~6 ist nicht eingelöst: gemeinsame Auswertung, keine
Kompatibilitätsidentität.} In Substep (xi-d) werden die beiden Seiten
ausdrücklich in eine Auswertung gebracht:

\begin{quote}\small
\glqq it is only by forming [a suitable positive tensor power of] the determinant of
the localizations of arithmetic vector bundles \dots\ and then applying the
[suitably normalized \dots] log-volume to various regions \dots\ that we are able
to obtain \emph{completely comparable objects}, namely,
$\mathbb{R}_{\le -|\log(\Theta)|} := \{\lambda \in \mathbb{R} \mid \lambda \le
-|\log(\Theta)|\} \subset \mathbb{R}$; $-|\log(q)| \in \mathbb{R}$\grqq
\hfill\anchor{IUT-III}{10376--10385}
\end{quote}

Zusammen mit der oben zitierten Log-Volumen-Deutung arithmetischer Grade existiert
damit eine numerische Auswertung $\nu_C$ im Container, und sie steht ausdrücklich
da, statt erschlossen werden zu müssen. So viel wird zugunsten der Quelle
protokolliert. Es ist aber nicht das, wonach Pflicht~6 fragt. Die Pflicht verlangt
die Kompatibilitätsidentität $\nu_C \circ i_q = \nu_q$ über die Identifikation der
beiden Container hinweg, und diese Identität steht nicht im Text: Der Block
\glqq gesucht, nicht gefunden\grqq{} weiter unten hält die Suche fest --- und hält fest, dass
das eine vorhandene Kompatibilitätslemma, die Log-Link-Kompatibilität der
Log-Volumina, Kompatibilität \emph{innerhalb} einer vertikalen Spalte sichert und
nicht über die Identifikation hinweg. Da eine gemeinsame Auswertung Pflicht~6 nie
für sich allein einlöst (Abschnitt~\ref{sec:obligations}), liefert R6a nicht.

\medskip\noindent
Version 0.1 dieser Arbeit codierte diese Hälfte umgekehrt, auf der damals
geltenden disjunktiven Fassung von Pflicht~6, und hielt fest, \glqq die Fehlstelle
sitzt bei Pflicht~7, nicht bei Pflicht~6\grqq. Die Zitate sind unverändert, und keine
Quelle wurde neu geprüft; die Codierung folgt dem gehärteten Wortlaut. Siehe
\glqq Änderungen in Version 0.2\grqq.

\paragraph{R6b --- Pflicht~7 wird gesetzt, nicht abgeleitet.} Substep (xi-e) fasst (SHE)
operativ als Nebenbedingung an die zulässigen Operationen:

\begin{quote}\small
\glqq `expressibility relative to the arithmetic holomorphic structure in the
1-column' [cf.\ (SHE)] amounts precisely to `\emph{expressibility via operations
that are valid/executable/well-defined even when subject to the condition that
the pilot-object log-volume associated to the input data prime-strip \dots\ be
equal to the fixed value} $-|\log(q)| \in \mathbb{R}$'.\grqq
\hfill\anchor{IUT-III}{10412--10420}
\end{quote}

und Substep (xi-f) zieht daraus die Membership:

\begin{quote}\small
\glqq the construction of the subset $\mathbb{R}_{\le -|\log(\Theta)|} \subset
\mathbb{R}$ of possible pilot-object log-volumes of output data \emph{is subject
to the condition} that this construction of output data possibilities
\emph{constitutes, in particular, a construction} [perhaps only up to some sort of
`approximation', as a result of various indeterminacies] \emph{of the
pilot-object log-volume of the input data prime-strip}, namely,
$-|\log(q)| \in \mathbb{R}$. The inclusion
$-|\log(q)| \in \mathbb{R}_{\le -|\log(\Theta)|}$, hence also the inequality
$-|\log(q)| \le -|\log(\Theta)| \in \mathbb{R}$ \dots\ \emph{then follows
formally}.\grqq{} \hfill\anchor{IUT-III}{10437--10451}
\end{quote}

Substep (xi-g) benennt das Ergebnis als \glqq two \emph{tautologically equivalent} ways
to compute the log-volume of the $q$-pilot object\grqq
\anchor{IUT-III}{10453--10458}, und die Rekapitulation des Autors wiederholt das
Muster: Die $q$-Verflechtung \glqq may be thought of as a special case of the
$\Theta$-intertwining \dots\ Corollary 3.12 then follows, essentially formally\grqq
\anchor{IUT-III}{10723--10731}.

\paragraph{Das Toy-Modell schreibt die tragende Annahme als Annahme aus.} Ohne
Labels widersprechen die beiden Zuordnungen einander
\anchor{IUT-III}{10747--10750}; der entscheidende Schritt wird dann als
Supposition eingeführt:

\begin{quote}\small
\glqq \emph{Now suppose that one verifies that one may construct} the `assignment with
indeterminacies' \dots\ as a mathematical structure that is intrinsically
associated to the underlying structure of the assignment $\Xi_q$ \dots\
\emph{even if one forgets the labels} `$q$', `$\Theta$' \dots, i.e., \emph{even if
one identifies} $q\mathbb{R}$, $\Theta\mathbb{R}$, in the usual way, with
$\mathbb{R}$ [cf.\ the properties (IPL), (SHE)]\grqq
\hfill\anchor{IUT-III}{10778--10784}
\end{quote}

woraus \glqq one then concludes formally\grqq{} die gewünschte Schranke folgt
\anchor{IUT-III}{10806--10813}.

\paragraph{Der Befund in Zertifikatssprache.} Im Realfall wird die Supposition
durch die Konjunktion (IPL)\,$\wedge$\,(SHE) eingelöst --- also durch
\emph{Eigenschaften, die dem Output eines Algorithmus zugesprochen werden}, nicht
durch ein Kompatibilitätslemma $\nu_C \circ i_q = \nu_q$. An derselben Stelle wird
zudem die von Pflicht~3 verlangte Container-Typisierung bewusst aufgegeben: Die
Identifikation von $q\mathbb{R}$ mit $\Theta\mathbb{R}$ ist Voraussetzung des
Schlusses, nicht sein Ergebnis. Das fehlende Element ist damit an beiden Hälften
der Zeile ein und dasselbe, und es fehlt doppelt: Die Identität ist das, was
Pflicht~6 verlangt, um die beiden Readouts durch Beweis statt durch Kohabitation
auf eine Skala zu bringen, und sie ist das, was Pflicht~7 bräuchte, um die
Ordnungsrelation über diese Skala hinweg zu tragen.

\paragraph{Fairness: \glqq tautologically\grqq{} ist affirmativ gemeint.} Die Quelle
verwendet dieses Wort zustimmend, nicht als Eingeständnis: Sie spricht davon,
\glqq nontrivial, albeit essentially tautological, consequences \dots\ such as the
inequality of Corollary 3.12\grqq{} zu ziehen \anchor{IUT-III}{8122--8125}. Für ihren
Autor ist die Tautologizität die Leistung, nicht der Mangel. Das Wort als
Selbstwiderlegung zu lesen wäre ein umgekehrtes Autoritätsargument und wird hier
nicht getan. Das Ledger fragt nur, ob die Kompatibilitätsidentität ausgewiesen und
ob die Ordnungsrelation \emph{abgeleitet} ist --- und nach dem Wortlaut ist
beides nicht der Fall.

\subsection{R8 --- ist ein zweiter Modus sichtbar oder ausgeschlossen?}

\textsc{Teilweise}, mit hoher Konfidenz. Benannt und qualitativ bepreist, aber
nicht als zweite Skala durch die Brücke geführt; auch nicht ausgeschlossen.

Die Doppelstruktur wird ausdrücklich genannt: eine Prime-Strip, die
\glqq \emph{simultaneously equipped with two intertwinings}\grqq{} ist, \glqq namely, the
$\Theta$-intertwining, \emph{up to indeterminacies}, and the $q$-intertwining\grqq
\anchor{IUT-III}{8114--8125}. Auch der Preis der Konstruktion wird deklariert:

\begin{quote}\small
\glqq this `solution' may be implemented \emph{only at the cost of admitting the
`indeterminacy' constituted by the upper semi-compatibility of} (Ind3).\grqq
\hfill\anchor{IUT-III}{10019--10031}
\end{quote}

und der Gehalt dieses Begriffs steht im Text selbst: \glqq upper semi-compatibility,
i.e., in a word, \emph{one-sided inclusions, as opposed to precise equalities}\grqq
\anchor{IUT-III}{8489--8490}.

Was nicht geschieht, ist die Führung des zweiten Modus durch den strittigen
Schritt. Vor der Brücke wird er in eine Richtung aufgelöst --- (Ind3) wird \glqq into
an inequality [from above]\grqq{} überführt ---, und an der Brücke wird die interne
Struktur durch (HIS) ausdrücklich vergessen. Sichtbar vor dem Schritt, suspendiert
in ihm.

Ein Beinahe-Treffer wird protokolliert statt unterschlagen: Eine Bemerkung, deren
Titel (\glqq The log-shifted nature of the localization homomorphism\grqq) eine sichtbar
geführte zweite Skala nahelegt, betrifft nur den Lokalisierungs-Homomorphismus
einer früheren Proposition \anchor{IUT-III}{5551--5553} und trägt die Zeile nicht.

\subsection{Ergebnis von Fall A}

Die resultierenden Befunde über die auditierten Zeilen sind in Tabelle~\ref{tab:case_a_findings} zusammengefasst.

\begin{table}[htbp]
\centering
\small
\renewcommand{\arraystretch}{1.12}
\begin{tabular}{@{}lP{0.60\textwidth}l@{}}
\toprule
Zeile & Befund & Konfidenz \\
\midrule
R2 & \textsc{liefert} --- Readout definitorisch fixiert und indeterminacy-frei;
  kein Transportschritt & hoch \\
R4 & \textsc{teilweise} --- Ort benannt, Last nicht beziffert, Bezifferung
  ausgelagert & hoch \\
R6a & \textsc{liefert nicht} --- gemeinsame Auswertung vorhanden und explizit;
  Kompatibilitätsidentität $\nu_C \circ i_q = \nu_q$ fehlt & hoch \\
R6b & \textsc{liefert nicht} --- Membership über ein Eigenschaftspaar gesetzt;
  folgt \glqq formally\grqq{} & hoch \\
R8 & \textsc{teilweise} --- benannt und bepreist, an der Brücke suspendiert & hoch \\
\bottomrule
\end{tabular}
\caption{Fall A: Zusammenfassung der Befunde über die auditierten BC-Ledger-Zeilen.}
\label{tab:case_a_findings}
\end{table}

Die Adresse der beiden nicht eingelösten Pflichten ist ein einzelner Substep. Dort
folgt die Membership \glqq then formally\grqq{} aus der Bedingung, dass die Konstruktion der
Output-Möglichkeiten \emph{zugleich} eine Konstruktion des Input-Readouts sei
\anchor{IUT-III}{10437--10451}. Was fehlt, ist eine Kompatibilitätsidentität
$\nu_C \circ i_q = \nu_q$ über die Container-Identifikation hinweg, die Pflicht~6
aus eigenem Recht verlangt, und darüber hinaus ein Vertragslemma, das die
Ordnungsrelation aus der erlaubten Linkage \emph{ableitet}, wie Pflicht~7 es
verlangt. An ihrer Stelle steht das Eigenschaftspaar, dessen Einlösung das
Toy-Modell als Supposition ausschreibt. Was der Text an diesem Substep sehr wohl
liefert, ist die gemeinsame Auswertung selbst, und zwar ausdrücklich --- das ist
nicht nichts; aber eine gemeinsame Auswertung ohne die Identität ist genau die
Kohabitation, die das Zertifikat von einer numerischen Aussage zu trennen gebaut
wurde.

\paragraph{Gesucht, nicht gefunden.} Über den gesamten Extrakt und getrennt über
den Beweisbereich: ein Kompatibilitätslemma der Form $\nu_C \circ i_q = \nu_q$
über die Container-Identifikation (Suchbegriffe: \texttt{compatib},
\texttt{log-volume}, \texttt{equal}, \texttt{identif}); ein Zahlwert für den
Approximationsterm (\texttt{approximation}, \texttt{epsilon}, \texttt{bound},
\texttt{estimate}, \texttt{Ind3}); eine explizite Verbuchung der
$O(\ell^2)$-Last an der Brücke (\texttt{theta value},
\texttt{procession-normalized}, \texttt{average}); eine sichtbar geführte zweite
Auswertungsskala an der Brücke (\texttt{log-shift}, \texttt{vertical shift},
\texttt{two intertwinings}). Eine wichtige Präzisierung, damit der Negativbefund
nicht als Nachlässigkeit gelesen wird: Ein echtes Kompatibilitätslemma
\emph{existiert} --- die Log-Link-Kompatibilität der Log-Volumina
\anchor{IUT-III}{10268--10272} ---, aber es liegt auf der falschen Achse. Es
sichert Kompatibilität \emph{innerhalb} einer vertikalen Spalte, nicht über die
Identifikation der beiden Auswertungscontainer hinweg --- und um genau diese geht
es in Pflicht~6, auf der Pflicht~7 dann erst aufbauen müsste.

\paragraph{Reichweite.} Hohe Konfidenz gilt den Textbefunden, Zeile für Zeile, im
autorseitigen Sinn, wie ihn Abschnitt~\ref{sec:instrument} festlegt. Die Reichweite
ist doppelt beschränkt: Geprüft wurde hier ausschließlich der dritte Band, und
gefüllt wurden nur vier der neun Zeilen --- berichtet wird also ein Befund an vier
Zeilen, kein Urteil des ganzen Zertifikats. Die von der Quelle selbst benannte
quantitative Fortsetzung ist Gegenstand von Fall B.

% =============================================================================
\section{Fall B: der deklarierte Quantifizierungsort}
\label{sec:caseB}
% =============================================================================

Fall A ließ eine konkrete Möglichkeit offen: Die an der Brücke nicht verbuchte
Last könnte an dem Ort verbucht sein, an den der Text sie ausdrücklich
auslagert. Die zweite Anwendung öffnet diesen Ort~\cite{Mochizuki2021d} und
stellt zwei strikt getrennte Fragen.

\begin{description}
\item[F1] Wird die $j^2$-/Prozessions-Last dort tatsächlich beziffert?
\item[F2] Leitet dieser Abschnitt die strittige Membership \emph{ab}, oder setzt
  er das Korollar \emph{voraus}?
\end{description}

Die Trennung ist wesentlich, denn eine positive Antwort auf F1 wird leicht als
Fortschritt bei F2 missverstanden, und das ist sie nicht.

\subsection{F1: beziffert}

\textsc{Beziffert}, mit hoher Konfidenz. Die Kette ist in der Quelle
ausgeschrieben und wird hier nur nachgewiesen, nicht nachgerechnet.

Der Theta-Wert geht als Eingangsgröße der Rechnung ein: \glqq $\lambda$\grqq{} wird als
\glqq $\ordd(-)$\grqq{} des Elements $q_{v_j}^{j^2}$ an den schlechten Stellen gesetzt
\anchor{IUT-IV}{1594--1597}. Der lokale Term trägt $j^2$ sichtbar, als
$j^2/(2\ell(j+1))$ \anchor{IUT-IV}{1656--1657}; nach Gewichtung bleibt
$j^2/(2\ell)$ \anchor{IUT-IV}{1669--1670}. Die Prozessions-Mittelung wertet dann
die Quadratsumme explizit aus:

\begin{quote}\small
\glqq it remains to compute the average over $j \in \mathbb{F}_\ell^\star$.
\emph{By averaging over} $j \in \{1,\dots,\ell^\star = (\ell-1)/2\}$ \emph{and
applying} (E1), (E2), we obtain the `procession-normalized upper bound'
\dots\ $= \frac{\ell+5}{4}\log(d_K) - \frac{\ell+1}{24}\log(q) + \dots$\grqq
\hfill\anchor{IUT-IV}{1675--1684}
\end{quote}

die dabei angewandte Identität lautet $\frac{1}{n}\sum_{m=1}^{n} m^2 =
\frac{1}{6}(2n+1)(n+1)$ \anchor{IUT-IV}{1379--1380}. Nach der lokalen
Abschätzung überlebt der Koeffizient bis in die Endschranke, wo er als
$-\frac{1}{6}\bigl(1 - \frac{12}{\ell^2}\bigr)\log(q)$ erscheint
\anchor{IUT-IV}{1775--1776}, und ebenso in der Aussage des Theorems selbst
\anchor{IUT-IV}{1352--1353}. Die verbleibende $\ell$-Abhängigkeit wird benannt und
als \glqq an inevitable consequence of the fundamental role played \dots\ by the
$\ell$-torsion points\grqq{} diskutiert \anchor{IUT-IV}{1909--1912}.

Die dargestellte Kette lässt keinen Term unverbucht: Die Last entsteht
nur an den schlechten Stellen, die guten nichtausgezeichneten Stellen liefern
null, der archimedische Schritt wird eigens gerechnet, und der Schlussschritt
summiert.

\subsection{F2: setzt voraus}

\textsc{Setzt voraus}, mit hoher Konfidenz, und der stärkste Beleg ist negativ.
Drei vollständige Suchen über das ganze Dokument liefern: null Treffer für den
$q$-Piloten; null Treffer für jede der vier deklarierten Eigenschaften, auf denen
die Brücke in Fall A ruht; null Treffer für die Verflechtungen. Alle Vorkommen
von \glqq pilot\grqq{} in diesem Dokument sind $\Theta$-Piloten. Ein Dokument, das unter den
gesuchten Begriffen den $q$-Piloten nie nennt und die tragenden Eigenschaften nie
erwähnt, liefert keine Herleitung der Brücke. Die nicht eingelöste Pflicht aus
Fall A liegt außerhalb des Gegenstands dieses Abschnitts, soweit diese Begriffe
reichen.

Die Rahmung sagt dasselbe positiv. Der Abschnitt wird eingeführt als Erzeugung
\glqq \emph{more explicit versions} \dots\ \emph{of the log-volume estimates for
$\Theta$-pilot objects obtained in} [IUT-III], Corollary 3.12\grqq
\anchor{IUT-IV}{456--458}, und sein Hauptsatz eröffnet mit \glqq \emph{Suppose that we
are in the situation of} [IUT-III], Corollary 3.12\grqq
\anchor{IUT-IV}{1281}.

Ebenso die Beweislogik. Der Abschnitt liefert die obere Seite --- \glqq we thus obtain
the following \emph{upper bound}\grqq{} \anchor{IUT-IV}{1773}, \glqq substituting back
\dots\ we obtain the following \emph{upper bound for} `$C$'\grqq
\anchor{IUT-IV}{1821} --- und importiert die untere:

\begin{quote}\small
\glqq and hence, \emph{by applying the inequality} `$C_\Theta \ge -1$' \emph{of} [IUT-III],
Corollary 3.12, conclude that \dots\grqq{} \hfill\anchor{IUT-IV}{1355--1356}
\end{quote}

Diese importierte Ungleichung \emph{ist} die strittige Membership in
Konstantenform: Die Konklusion des Korollars lautet
$-|\log(\Theta)| \ge -|\log(q)|$, \glqq i.e., $C_\Theta \ge -1$ for any real number $C_\Theta$
such that $-|\log(\Theta)| \le C_\Theta \cdot |\log(q)|$\grqq{} \anchor{IUT-III}{9878--9884}. Die
Architektur, die der Text sich selbst aufschreibt, lautet also: Dieser Abschnitt
begrenzt $C$ nach oben, das Korollar begrenzt $C$ nach unten, und erst beides
zusammen sagt etwas über $\log(q)$. Eine obere Schranke für $C$ allein sagt
nichts.

Schließlich behandelt der Abschnitt nicht einmal die dritte Indeterminacy selbst.
An allen drei Stellen, an denen sie auftritt, wird die Behandlung zurückdelegiert:
Sie werde \glqq taken into account by the fact that we are considering upper bounds
[cf.\ the discussion of Step (x) of the proof of \dots\ Corollary 3.12]\grqq
\anchor{IUT-IV}{1609--1611}.

\subsection{Konsequenz für das Ledger}

Drei Konsequenzen, und nur die erste ist eine Verdiktänderung.

\begin{enumerate}
\item \textbf{Zeile R4 spaltet sich auf.} Die Last wird in der Auswertung der
  $\Theta$-Seite beziffert (R4a: ja, im Verbund der beiden Bände) und nicht an
  der Brücke verbucht (R4b: nein im dritten Band, nicht Gegenstand im vierten).
  Ohne die Trennung verdeckt ein Wort zwei Befunde.
\item \textbf{Die Auslagerung ist verifiziert.} Die Sorge, die die Zeile
  motivierte --- eine Last, die in einer Koordinate verschwindet und unverbucht
  zurückkehrt ---, trifft für die beiden Bände zusammen nicht zu: Der Theta-Wert
  verschwindet im dritten und \emph{kehrt im vierten benannt und gerechnet
  wieder}, an genau der Stelle, die der dritte angegeben hatte. Das wird zugunsten
  der Quelle festgehalten.
\item \textbf{Das R6-Verdikt bleibt an beiden Hälften unverändert, seine
  Reichweite wächst.} Fall A hatte sein Verdikt unter den Vorbehalt gestellt,
  dass nur ein Band geprüft wurde. Für den von diesem Band benannten Ort ist der
  Vorbehalt jetzt abgearbeitet: Dort steht weder eine Kompatibilitätsidentität
  noch ein Vertragslemma, und kein Kandidat für eines von beiden. Der Ort ist
  nicht widersprüchlich, sondern thematisch unbeteiligt.
\end{enumerate}

Eine weitere Beobachtung gehört hierher, als Beobachtung zur Architektur und nicht
als Urteil: Indem der vierte Band $C_\Theta \ge -1$ als Eingabe nimmt, macht er sichtbar,
wie viel an der Membership hängt. Seine gesamte explizite Rechnung produziert eine
obere Schranke, die ohne die vom Korollar gelieferte untere Schranke gar keine
Aussage über $\log(q)$ ergibt.

\paragraph{Gesucht, nicht gefunden.} Eine Ableitung der Membership irgendwo in
diesem Band (Suchbegriffe: \texttt{q-pilot}, \texttt{pilot}, \texttt{holomorphic
hull}, \texttt{inclusion}, \texttt{follows formally}, \texttt{tautolog},
\texttt{Step (xi)}): nicht gefunden; der eine Treffer zum Substep liegt in einem
späteren Abschnitt und verweist auf den dritten Band zurück. Ein Vertragslemma
oder eine Diskussion des Eigenschaftspaars: nicht gefunden, null Treffer. Ein
Zahlwert für den Approximationsterm des Toy-Modells: nicht gefunden --- und die
naheliegende Gleichsetzung dieses Terms mit der hier gerechneten Schranke ist
\emph{unsere} Zuordnung, kein Textbefund; der Text stellt sie nirgends her.

\paragraph{Import-Ledger.} Das Verdikt von F2 hängt an einem einzigen Rückgriff auf
fremdes Material, und die Disziplin aus Abschnitt~\ref{sec:importledger} verlangt,
ihn einzutragen statt zu übergehen (Tabelle~\ref{tab:case_b_import_ledger}).

\begin{table}[htbp]
\centering
\small
\renewcommand{\arraystretch}{1.12}
\begin{tabular}{@{}P{0.24\textwidth}P{0.66\textwidth}@{}}
\toprule
Feld & Eintrag \\
\midrule
Importierende Stelle & IUT-IV, \S\,1, im Beweis der dortigen Hauptabschätzung
  \anchor{IUT-IV}{1355--1356} \\
Importierte Aussage & die Ungleichung \glqq$C_\Theta \ge -1$\grqq{} aus IUT-III, Corollary 3.12 \\
Benötigte Aussage & dieselbe Ungleichung: Der Abschnitt erzeugt eine
  obere Schranke für $C$ und braucht eine untere, die er nicht herleitet \\
Passung & exakt --- importierte und benötigte Aussage fallen zusammen \\
Geprüft wurde & die Passung, nicht der Beweis, gemäß der asymmetrischen Regel aus
  Abschnitt~\ref{sec:importledger} \\
\bottomrule
\end{tabular}
\caption{Fall B: Import-Ledger-Eintrag für die IUT-IV-Hauptabschätzung.}
\label{tab:case_b_import_ledger}
\end{table}

\noindent
Ein Z-Typ-Versagen wird hier nicht protokolliert: Der Import trägt genau die
Aussage, für die er aufgerufen wird. Gerade deshalb zählt der Eintrag. Die
importierte Aussage ist diejenige, nach deren Herleitung Pflicht~7 fragt --- der
Rückgriff ist also einwandfrei, und der Ort kann die Pflicht trotzdem nicht
einlösen.

% =============================================================================
\section{Fall C: eine vergleichende Diskriminierungsanwendung in einer anderen
Theorie}
\label{sec:caseC}
% =============================================================================

Die dritte Anwendung ist kein zweiter Test desselben Textes. Sie wendet das
unveränderte Instrument auf ein unabhängiges Vergleichsprogramm an~\cite{JoshiATSIII, JoshiATSIV}, das dieselbe Ungleichung auf anderem Weg
erreicht. Ihre Funktion ist die Beantwortung einer Frage, die eine einzelne
Anwendung nicht beantworten kann: Liefert das Zertifikat überall denselben Befund
--- dann ist es eine Schablone und kein Instrument --- oder unterscheidet es?

Das Etikett ist von Bedeutung und wird hier in den Begriffen des
Methodenpapiers~\cite{GeigerBCM2026} festgelegt: Dies ist eine \emph{vergleichende
Diskriminierungsanwendung}, keine Kontrolle. Version 0.1 nannte sie
Vergleichskontrolle, was mehr beanspruchte, als der Fall tragen kann --- eine
Kontrolle müsste ein Transportargument sein, dessen siebte Pflicht \emph{eingelöst
ist}, und Abschnitt~\ref{sec:convergence} hielt bereits fest, dass kein solcher
Fall unter den geprüften war. Die unten beschriebene Funktion bleibt unverändert;
korrigiert wird allein der Name.

Es unterscheidet. Drei der vier Zeilen liefern, die vierte liefert nur
eingeschränkt, und die nicht eingelöste Pflicht sitzt an anderer Adresse und hat
eine andere Gestalt. \emph{Wo} diese nicht eingelöste Pflicht sitzt, haben zwei
spätere Korrekturen revidiert (Abschnitt~\ref{sec:caseC-nachtraege}); der
Vergleich am Ende dieses Abschnitts gibt die revidierte Lesart wieder, nicht die
erste Codierung, und die Revision wird als Teil des Falls berichtet.

\subsection{R2: geliefert, über eine bewiesene Relation}

Der Transport entlang des konstruktionseigenen Links erhält den Readout, und der
$j^2$-Faktor wird sichtbar mitgeführt: Die Skalenrelation der Bewertungen steht
als Teil eines Satzes, $|-|_{K'_{y_{w,j}}} = |-|^{j^2}_{K'_{y_{w,1}}}$ für
$j = 1,\dots,\ell^*$ \anchor{ATS-III}{1544--1560}, und wird in der expliziten
Berechnung der Bloch--Kato-Klassen der Theta-Werte verwendet
\anchor{ATS-III}{6458--6484}. Die Auswertungsskala ist konkret: ein Log-Volumen
mit angegebener Normierung \anchor{ATS-III}{6562--6596}. Pflicht~1 ist eingelöst;
Pflicht~4 wird hier durch eine \emph{bewiesene Bewertungsrelation} eingelöst,
nicht durch eine Stipulation.

Das ist der Kontrast, der die Aufteilung R2a/R2b erzwungen hat. Es tritt auch ein
zweiter Transport auf, zwischen zwei durch einen globalen Frobenius verbundenen
Kopien der Umgebungsdaten, und dort wird die Identifikation \emph{durch eine
Normierungswahl gesetzt}: \glqq Choose the same global normalization \dots\ With this
choice \dots\ one has\grqq{} --- dies in der \emph{Aussage} von Theorem 6.10.1 aus
ATS-IV \anchor{ATS-IV}{3186--3189}; der Beweis der resultierenden Gleichheit ist,
im Beweis eben dieses Satzes, der eine Satz, sie sei \glqq a consequence of the choice
of the same global normalization\grqq{} \anchor{ATS-IV}{3262--3266}. Dieselbe Zeile,
zwei verschiedene Arten von Antwort.

\subsection{R4: geliefert, über einen globalen Zwang}

Die Last entsteht nur an den semistabilen Stellen, überlebt bis in die numerische
Auswertung, und ihre Kompensation wird verbucht --- Buchungsinstrument ist die
Produktformel:

\begin{quote}\small
\glqq the non-trivial scaling relationship \dots\ at primes $w \in
V^{\mathrm{odd},ss}$ \emph{forces that suitable (re)normalization of valuations
must be introduced} at primes $w \in V_{L'} - V^{\mathrm{odd},ss}$ so that the
product formula \dots\ hold[s] for each normalized arithmeticoid\grqq
\hfill\anchor{ATS-III}{1600--1609}
\end{quote}

Ehrlichkeit verlangt die Einschränkung daneben: Der Beweis dieser Aussage ist ein
Satz, weshalb die \emph{Existenz} einer konsistenten kompensierenden Renormierung
aus dem Zwang erschlossen und nicht konstruiert wird. Die Last ist verortet, nicht
beziffert. Der Fehlermodus, dessentwegen die Zeile existiert, tritt nicht ein.

\subsection{R6: teilweise --- als Ungleichung bewiesen, als Brücke konventionell}

Die Ungleichung selbst ist ein Satz mit Beweis --- Theorem 9.11.1 aus ATS-III, in
dessen \S\,9.11 \anchor{ATS-III}{6794--6806} --- und
die Auswertung im Container ist echt numerisch, berechnet als Summe lokaler
Beiträge \anchor{ATS-IV}{764--769}. Pflicht~6 ist \emph{teilweise} eingelöst: Die
Auswertung ist numerisch, aber die Kompatibilität hält nur zwischen den Summen und
wird stellenweise von der Quelle selbst dementiert --- wie die im nächsten Absatz
zitierten Passagen es in den Worten der Quelle festhalten. Version 0.1 codierte
diese Hälfte als \glqq über ihre erste Alternative\grqq{} eingelöst; unter dem gehärteten
Wortlaut von Pflicht~6 ergeben dieselben Zitate den teilweisen Code, und die
Begründung stand bereits im Text (\glqq Änderungen in Version 0.2\grqq).

Die Brücke, die den Schluss trägt, ist eine andere Sache. Zwei Auswertungen in
zwei durch den globalen Frobenius verbundenen Kopien werden durch eine mittlere
Gleichheit verbunden, deren Beweis im Beweis von Theorem 6.10.1 vollständig
lautet:

\begin{quote}\small
\glqq The middle equality is a tautology using the fact that the number fields
provided by $y_0'$ and $y_0$ are identically normalized.\grqq
\hfill\anchor{ATS-IV}{3272--3273}
\end{quote}

während derselbe Text den stellenweisen Vergleich, mit dem ein Leser sie prüfen
könnte, ausdrücklich verbietet: \glqq one cannot claim that $z_p \le z_p'$ holds for
all primes $p$ without violating the middle equality\grqq
\anchor{ATS-IV}{816--819}, und der Unterschied durch Frobenius an jeder Stelle
\glqq \emph{precludes direct comparison of various local quantities}\grqq
\anchor{ATS-IV}{3208--3211}. Ordnungskompatibilität besteht damit nur zwischen den
Summen, nie stellenweise.

\subsection{R8: geliefert --- und deshalb ist der Rest lesbar}

Der zweite Modus ist hier der globale Frobenius-Shift, und er ist nicht bloß
vorhanden, sondern deklariert, asymmetrisch eingesetzt, gedeutet und als Korrektur
der eigenen früheren Fassung markiert: Die untere Schranke wird in der einen, die
obere in der anderen Kopie berechnet, \glqq the upper bound is not independent of this
choice\grqq{} \anchor{ATS-IV}{733--743}, und eine eingeklammerte Notiz hält fest, dass
\glqq [t]he previous release of this paper had omitted this Frobenius-shifting\grqq
\anchor{ATS-IV}{755}.

Das ist das stärkste Positivergebnis des Vergleichs, und sein Ertrag ist
analytisch, nicht rhetorisch: \emph{Weil} der zweite Modus offengelegt statt in
einer Normierung absorbiert wird, hat die Lücke bei R6 überhaupt eine Adresse.

\subsection{Zwei spätere Korrekturen, die in entgegengesetzte Richtungen ziehen}
\label{sec:caseC-nachtraege}

Beide sind zu berichten, sonst wird der Fall falsch dargestellt.

\paragraph{Stärkung.} Eine Quellenprüfung der beiden früheren Bände derselben
Serie, gesondert als Werkzeug-Supplement dieser Serie
veröffentlicht~\cite{GeigerGNC2026}, hat die oben offene Definitionsfrage
entschieden: Das normierte Objekt ist pro Punkt eindeutig, \glqq the same global
normalization\grqq{} ist also \emph{definitorisch}. Auf dieser Lesart ist die Identifikation nicht lasttragend,
und die mittlere Gleichheit ruht stattdessen auf der konstruktionsbedingten
Frobenius-Stabilität der kollationierten adelischen Menge --- also genau auf dem
Mechanismus, den diese Anwendung unter R8 als geliefert eingestuft hat --- plus
einer definitorischen Skalenpinnung. Übrig bleibt kein fehlendes Vertragslemma,
sondern eine engere Frage nach einer Maßkonvention.

\paragraph{Schwächung.} Eine Substanzprüfung des R6-tragenden Satzes ---
Begleitarbeit dieser Serie, erschienen als Teil~5~\cite{GeigerSubstance2026} --- hat
ergeben, dass sein Kernschritt, ein als \glqq clear from the construction\grqq{} deklariertes
Enthaltensein, von der dafür zitierten Konstruktion nicht getragen wird. Das
ist eine G+D-Lücke: ein fehlendes Argument (eine Supremum-Aussage steht, wo eine
Maß-Aussage gebraucht wird) zusammen mit einer strukturellen Unschärfe darüber, wo
die eine algebraische Struktur zur anderen wird. Die Reparaturpflicht ist benannt;
der Befund ist kein Fehler-Nachweis.

Beides zusammen fasst sich am besten in der Formulierung dieser beiden
Prüfungen: Der Unterschied in der \emph{Form} bleibt, der Vorsprung in der
\emph{Substanz} schrumpft.

\paragraph{Die revidierte Codierung.} Zusammengenommen machen die beiden
Korrekturen diesen Fall weder besser noch schlechter als die zuerst festgehaltene
Codierung; sie \emph{verschieben} das Defizit. Die Zeilen R2, R4 und R8 sind von
keiner der beiden Korrekturen berührt und liefern weiterhin. Was sich ändert, ist
die Aufteilung innerhalb des \textsc{teilweise} von Zeile R6. Vor den Korrekturen
war der eingelöste Teil die Ungleichung als bewiesener Satz und der fehlende Teil
eine Brücke, die über eine konventionelle Identität geschlossen wird. Nach ihnen
ist die Identität definitorisch und damit nicht der fehlende Teil, während der
fehlende Teil im Satz selbst sitzt, an dessen Enthaltensein-Schritt. Das Verdikt
bleibt \textsc{teilweise}, weil die Regel lautet, dass ein geteiltes Verdikt seine
Aufteilung nennen und nicht mitteln muss --- und geändert hat sich die Aufteilung,
nicht das Verdikt.

Dass eine spätere, tiefere Lesart die erste eigene Codierung revidiert hat, wird
hier berichtet und nicht geglättet. Ein Eintrag, den eine erneute Prüfung
derselben Quellen korrigieren kann, ist ein prüfbarer Eintrag; das ist eine
Eigenschaft, die man haben will. Es ist keine Behauptung, das Instrument sei
validiert.

\subsection{Was der Vergleich zeigt}

Beide Argumente verfehlen Pflicht~7 im strengen Sinn des Ledgers, und sie
verfehlen sie unterschiedlich --- was der ganze Sinn einer vergleichenden
Diskriminierungsanwendung ist (Tabelle~\ref{tab:case_a_c_comparison}).

\begin{table}[htbp]
\centering
\small
\renewcommand{\arraystretch}{1.12}
\begin{tabular}{@{}lP{0.40\textwidth}P{0.40\textwidth}@{}}
\toprule
Zeile & Fall A & Fall C \\
\midrule
R2 & liefert --- kein Transport (deklariert) & liefert --- bewiesene
  Bewertungsrelation; zweiter Transport konventionell \\
R4 & teilweise --- Ort benannt, Bezifferung ausgelagert & liefert --- verbucht
  über den Produktformel-Zwang \\
R6 & liefert nicht an beiden Hälften --- keine Kompatibilitätsidentität, und
  Membership folgt \glqq formally\grqq{} & teilweise, revidiert --- Kompatibilität nur
  zwischen den Summen, stellenweise dementiert; Ungleichung von einem Satz
  getragen, Defizit im Enthaltensein-Schritt dieses Satzes \\
R8 & teilweise --- benannt, dann suspendiert & liefert --- offengelegt und
  asymmetrisch genutzt \\
\bottomrule
\end{tabular}
\caption{Komparative Diskriminationsmatrix: BC-Ledger-Befunde in Fall A (IUT-III) gegenüber Fall C (ATS-IV).}
\label{tab:case_a_c_comparison}
\end{table}

Der Unterschied, auf den es ankommt, ist die \emph{Gestalt} der nicht eingelösten
Pflicht, und die Korrekturen oben haben verschoben, wo diese Gestalt in Fall C
sitzt, ohne zu ändern, dass es ein identifizierbarer Schritt ist. In Fall C ist das
Defizit lokalisiert: Nach den Korrekturen sitzt es an einem einzelnen
Enthaltensein-Schritt innerhalb des Satzes, der die Ungleichung trägt, während es
zuvor an der konventionellen Identität gelesen wurde, die die beiden Auswertungen
verbindet. In Fall A ist sie flächig: eine Konjunktion zweier deklarierter
Eigenschaften, verteilt über eine Bemerkung und vier Substeps, eingelöst durch
einen Schluss, den der Text selbst formal nennt. Der Gegensatz ist der zwischen
einem Defizit an einem identifizierbaren Schritt und einem über mehrere verteilten.
Er ist keine Aussage darüber, was schwerer wiegt --- die Typologie aus
Abschnitt~\ref{sec:typology} typisiert Lücken und graduiert sie bewusst nicht.

% =============================================================================
\section{Fall D: nachgelagerte Anwendungen}
\label{sec:caseD}
% =============================================================================

Die vierte Anwendung stellt zwei Arbeiten, die die strittige Ungleichung
\emph{benutzen}~\cite{ZhouI, ZhouII}, eine andere Frage: Setzen sie das
Zertifikat voraus, umgehen sie es, oder führen sie eine alternative Skala ein? Ein
Befund hier kann die Lücke nicht detektieren --- er kann nur zeigen, wie sie sich
fortpflanzt.

\paragraph{Anmerkung zum Etikett.} Diese Arbeiten sind im zugrundeliegenden
Material unter einem Namen abgelegt, der in ihnen nicht vorkommt; die verwendeten
Kurven sind Frey--Hellegouarch-Kurven, und der Autor ist ein einzelner Autor. Das
Etikett ist ein Ablage-Artefakt und wird hier nicht verwendet.

\subsection{Wo die Schranke eintritt}

In der ersten Arbeit erscheint die strittige Ungleichung als Proposition~1.9, deren
Formel~(1.7), und ihr gesamter Beweis ist ein Satz: Sie folge \glqq from the
$\mu_6$-version of\grqq{} dem Korollar aus Fall A \anchor{Zhou-I}{728--732}. Die
Formelnummer wird danach genau einmal verwendet, im Beweis der folgenden
Proposition \anchor{Zhou-I}{753--757}, und Proposition~1.9 wird nie wieder
aufgegriffen. Es gibt keine Herleitung, kein Zwischenlemma und kein Vertragslemma
der von Pflicht~7 verlangten Art.

Alles, was die Arbeit eigenständig beiträgt, liegt stromabwärts auf der
$\Theta$-Seite: eine schärfere Hüllen-Schranke \anchor{Zhou-I}{482--496}, ein
eigener Readout des $\Theta$-Piloten \anchor{Zhou-I}{677--725} und eigene
Konstanten in der Zielungleichung \anchor{Zhou-I}{742--751}, die die Arbeit selbst
als Verbesserung gegenüber einer früheren expliziten Abschätzung ausweist
\anchor{Zhou-I}{1931--1934}.

Die zweite Arbeit nennt das Korollar überhaupt nicht. Sie erbt die Brücke über die
erste Arbeit, und zwar genau an der Brückenproposition
\anchor{Zhou-II}{291--292}, und sie geht weiter: Eine Bemerkung verfolgt, wo eine
Hypothese in den vier zitierten Bänden verwendet wird, und schließt aus dieser
Verfolgung, deren Resultate gälten \glqq still \dots\ for\grqq{} eine andere Klasse von
Ausgangsdaten \anchor{Zhou-II}{267--280}. Der strittige Substep wird nie berührt.

\subsection{Schweigen, und ein aufschlussreicher Kontrast}

Suchen über beide Dokumente liefern null Treffer für die bekannte Kritik, für die
Indeterminacies, für die multiradiale Konstruktion, für die Piloten und für jeden
gesuchten Begriff aus dem Feld Streit oder Korrektheit. Die Positionierung
geschieht stattdessen implizit, durch das Wort \glqq unconditionally\grqq{} für die
Konsequenzen \anchor{Zhou-I}{29--30}.

Was das über eine bloße Abwesenheit hinaushebt, ist der Kontrast. Beide Arbeiten
praktizieren durchaus explizite Rigorositätsvorbehalte --- die zweite sagt offen,
dass sie einen Schritt auslässt, und \glqq claim[s] that we can always do the similar
things presented above\grqq{} \anchor{Zhou-II}{153--156} ---, wenden sie aber nur auf
eigene Auslassungen an, nie auf den importierten Schritt.

\subsection{Keine alternative Skala}

Die Log-Volumen-Normierung wird aus dem vierten Band der Quelltheorie zitiert
statt ersetzt \anchor{Zhou-I}{474--480}. Die schärferen Schranken, der eigene
Readout und die neuen Konstanten liegen sämtlich auf den Seiten der Brücke, nicht
in ihr. Eine Frage bleibt ausdrücklich unentschieden: ob ein expliziter Faktor in
der importierten Ungleichung eine Renormierung derselben Größe ist, die das
Korollar in seine Notation einfaltet, oder eine davon verschiedene Größe, war im
Quellenumfang jener Anwendung nicht entscheidbar und ist als Folgeprüfung vermerkt
statt beantwortet.

\subsection{Ergebnis von Fall D}

Beide Arbeiten \emph{setzen} das Zertifikat \emph{voraus}. Keine umgeht es, und
keine führt eine alternative Skala dafür ein. Die erste steuert ein Element von
Pflicht~6 bei --- ihren eigenen expliziten numerischen Readout auf der
$\Theta$-Seite ---, liefert aber keine Kompatibilitätsidentität, und sie
übernimmt Pflicht~7 vollständig. Die nicht eingelöste Pflicht wird also
stromabwärts nicht geschlossen, sondern unbesehen geerbt und in der zweiten Arbeit
zusätzlich auf eine Datenklasse ausgedehnt, die die zitierten Bände nicht
behandeln.

\paragraph{Gesucht, nicht gefunden.} Beide Extrakte wurden case-insensitiv
durchsucht; die Zahlen stehen als erste Arbeit\,/\,zweite Arbeit.
\texttt{Scholze} 0/0; \texttt{Stix} 0/0; \texttt{indetermin} 0/0;
\texttt{multiradial} 0/0; \texttt{Theta-pilot}, \texttt{q-pilot},
\texttt{theta pilot} 0/0; \texttt{Frobenioid} 0/0; \texttt{controvers},
\texttt{criticis}, \texttt{criticiz}, \texttt{dispute}, \texttt{debate} 0/0;
\texttt{gap in}, \texttt{not verif}, \texttt{unverif}, \texttt{validity},
\texttt{correctness} 0/0. Zwei Suchbegriffe liefern Treffer, und keiner davon ist
eine Auseinandersetzung: das maskierte \texttt{3.12} liefert 1/0, wobei der eine
Treffer die Zitation bei \anchor{Zhou-I}{732} ist, und \texttt{log-theta-lattice}
liefert 1/1, in beiden Fällen der Titel eines Literaturverzeichnis-Eintrags
\anchor{Zhou-I}{2583--2584}, \anchor{Zhou-II}{2254--2255}. Die Reichweite dieses
Blocks ist die Reichweite dieser Wortliste: Eine unter anderer Terminologie
geführte Auseinandersetzung wäre damit nicht erfasst worden.

\paragraph{Versionsvorbehalt.} Die zweite Arbeit bezeichnet sich in der ersten
Zeile ihres Abstracts selbst als ältere, noch zu aktualisierende Fassung
\anchor{Zhou-II}{10}. Jeder Befund über sie ist ein Befund über diese Fassung;
eine spätere Fassung könnte insbesondere den Vererbungsschritt anders begründen.

% =============================================================================
\section{Was die vier Anwendungen gemeinsam zeigen}
\label{sec:convergence}
% =============================================================================

\subsection{Vier verschiedene Rollen, nicht vier Detektionen}

Es wäre bequem und falsch, zu sagen, der Befund sei viermal bestätigt worden. Die
vier Anwendungen leisten Verschiedenes, und sie getrennt zu benennen ist zugleich
genauer und informativer (Tabelle~\ref{tab:applications_roles}).

\begin{table}[htbp]
\centering
\small
\renewcommand{\arraystretch}{1.12}
\begin{tabular}{@{}lP{0.20\textwidth}P{0.48\textwidth}@{}}
\toprule
Fall & Rolle & Was er feststellt \\
\midrule
A & \textbf{Detektion} & Die Pflichten~6 und~7 sind an einem benannten Substep
  nicht eingelöst: Die Kompatibilitätsidentität ist nicht ausgewiesen, und die
  Ordnungsrelation ist nicht abgeleitet. \\
B & \textbf{Weg\-schließung} & Der Ort, an den die Bezifferung ausgelagert wird,
  ist explizit und an den von ihm benannten Knoten geschlossen und kann Pflicht~7
  nicht einlösen, weil er die strittige Ungleichung als Eingabe verbraucht. \\
C & \textbf{Vergleichende Diskri\-minierung} & Das Instrument unterscheidet: drei Zeilen
  liefern, und die nicht eingelöste Pflicht sitzt an anderer Adresse, in anderer
  Theorie, mit anderer Gestalt --- an einer Adresse, die im Lauf dieser Arbeit
  durch eine genauere Lesart derselben Quellen revidiert wurde. \\
D & \textbf{Fort\-pflanzung} & Nachgelagerte Arbeiten erben die nicht eingelöste
  Pflicht ungeprüft und dehnen das Geerbte in einem Fall aus. \\
\bottomrule
\end{tabular}
\caption{Systematische Rollen und methodische Funktionen der vier Fallanwendungen.}
\label{tab:applications_roles}
\end{table}

Nur A ist eine Detektion. B ist das Gegenteil einer zweiten Detektion: Es stellt
fest, dass ein bestimmter Kandidatenort \emph{nicht} der Ort ist, an dem die
Pflicht eingelöst wird, was die Suche verengt statt sie zu wiederholen. C ist eine
Diskriminierungsanwendung, deren Wert darin liegt, eine \emph{andere} Antwort zu
liefern. D sagt
nichts darüber, ob die Pflicht eingelöst werden kann; es sagt, dass keine der
beiden geprüften Fassungen das tut.

Fall B trägt eine zweite Funktion, die zu benennen sich lohnt, weil sie der
einzige interne Beleg dafür ist, dass das Instrument auch etwas anderes als ein
Defizit zurückgeben kann. Unter denselben Kriterien, die in Fall A eine nicht
eingelöste Pflicht protokollieren, liefert der in Fall B geprüfte Text eine
explizite bezifferte Kette und einen exakt passenden Import an den tragenden
Knoten: Auf diese Fragen fällt die Antwort positiv aus. Das ist ein
\emph{positiver dokumentarischer Vergleich}, keine Validierung --- Fall B ist kein
Ort, an dem Pflicht~7 eingelöst werden könnte, und kann deshalb genau die Zeile
nicht prüfen, an der das Zertifikat dem Vorwurf am stärksten ausgesetzt ist, es sei
auf seine Antwort hin gebaut.

\subsection{Die Adresse ist über unabhängige Texte hinweg stabil}

Für eine vollständige synoptische Gesamtschau fasst Tabelle~\ref{tab:synoptic_cross_case_matrix} alle vier Fallanwendungen bezüglich Primärquellen, Theorierahmen, BC-Ledger-Kernbefunden, Lückentypen und methodischen Funktionen im Argumentationsgang zusammen.

\begin{table}[htbp]
\centering
\small
\setlength{\tabcolsep}{2.5pt}
\renewcommand{\arraystretch}{1.15}
\begin{tabular}{@{}lP{0.15\textwidth}P{0.21\textwidth}P{0.20\textwidth}lP{0.14\textwidth}@{}}
\toprule
Fall & Primärquelle & Theoretischer Kontext & BC-Ledger-Befund & Defizit & Rolle \\
\midrule
A & IUT-III, Kor.~3.12 & Inter-universelle Teichmüller-Theorie & R6a/R6b \textsc{liefert nicht} & G (Schritt xi-e) & Detektion \\
B & IUT-IV, \S\,1 & Explizite Log-Volumen-Quantifizierung & R4 \textsc{teilweise} / Import exakt & Z-frei & Weg\-schlie\-ßung \\
C & ATS-IV v2, \S\,6.10, Thm.~6.10.1 & Arithmetische Teichmüller-Räume & R6 \textsc{teilweise} (revidiert) & G (Summenebene) & Diskrim. \\
D & Zhou-I/II & Nachgelagerte diophantische Anw. & R6 \textsc{geerbt} & G (importiert) & Vererbung \\
\bottomrule
\end{tabular}
\caption{Synoptische Gesamtschau aller vier Anwendungen: Gegenstand, Kernbefunde, Lückentypen und methodische Rollen.}
\label{tab:synoptic_cross_case_matrix}
\end{table}

Was die vier gemeinsam feststellen, ist enger und, wie wir meinen, dauerhafter als
eine Zählung von Bestätigungen. \emph{Zwei} textlich unabhängige Werkkomplexe ---
die Theorie der Fälle~A und~B und das Programm aus Fall~C ---, mit einem
unveränderten Instrument geprüft, legen die Last des Arguments auf dieselbe
Passage: den Schritt von einem gemeinsamen Container zu einer abgeleiteten
Ordnungsrelation. Keiner von beiden vollzieht diesen Schritt, und keiner löst auch
die unmittelbar davorstehende Pflicht ein --- wobei sie dort unterschiedlich
ausfallen, was den Vergleich erst informativ macht. Ein dritter Komplex,
Fall~D, legt nichts: Er importiert den Schritt aus dem ersten, wie sein eigener
Text sagt, und kann deshalb nichts bestätigen --- was er beiträgt, ist die
Beobachtung, dass die Pflicht ungeprüft weiterwandert. Pflicht~6 ist in Fall~A
nicht eingelöst, weil die Kompatibilitätsidentität schlicht fehlt; in Fall~C
teilweise, weil die Kompatibilität zwischen den Summen hält und stellenweise
dementiert ist; und in Fall~D nur auf einer Seite ergänzt. Pflicht~7 ist im
strengen Sinn des Ledgers in keinem von ihnen eingelöst.

Dass das Instrument die Fälle dennoch scharf trennt --- \glqq liefert nicht\grqq{} gegen
\glqq teilweise\grqq{} gegen \glqq geerbt\grqq{} ---, spricht gegen die Lesart, der gemeinsame Befund sei
bloß ein Artefakt der Zertifikatskonstruktion. Entschieden ist die Frage damit
nicht. Eine echte Positivkontrolle --- ein Transportargument derselben Gestalt,
dessen siebte Pflicht \emph{eingelöst ist} --- war unter den geprüften Fällen nicht
dabei, und diese Lücke zu benennen gehört zum Befund und ist keine Randbemerkung.

\subsection{Ein unabhängiges Formalisierungsprojekt erreicht denselben Ort}

Im selben Zeitraum hat ein internationales Formalisierungsprojekt einen
Zwischenbericht veröffentlicht~\cite{LANA2026}, der an derselben Passage ein
Kompatibilitätsproblem isoliert. In seinen eigenen Worten:

\begin{quote}\small
\glqq we isolate a specific compatibility problem at the final stage of the argument:
the relation between the construction arising directly from the $q$-pilot in its
native arithmetic holomorphic structure and the construction obtained by anabelian
and Kummer-theoretic methods through the multiradial procedure. \dots\
\emph{Project LANA has not yet reconstructed a proof of the required
compatibility}.\grqq{} \hfill\anchor{LANA}{11--16}
\end{quote}

und zum Status dieser Aussage:

\begin{quote}\small
\glqq the proof of the final numerical inequality \emph{hinges on the compatibility}
(9-1) \emph{which is not manifestly false. However, we \dots\ do not have a proof
of} (9-1) \emph{at this time}.\grqq{} \hfill\anchor{LANA}{2988--2990}
\end{quote}

Die verlangte Aussage ist eine Existenzbehauptung: Konstruiert werden zwei
Isomorphismen punktierter reeller Vektorräume mit gemeinsamem Ziel --- der eine
aus dem $q$-Piloten \glqq together with the native arithmetic holomorphic structure\grqq
\anchor{LANA}{2807--2812}, der andere anabelisch und abhängig von der Wahl einer
Integralstruktur, \glqq subject to indeterminacies\grqq{} \anchor{LANA}{2783--2803} ---, und
gewünscht wird, \glqq that \emph{there exists some suitable} $S$ \emph{such that}
$\eta^q = \eta_S^{\mathrm{anab}}$\grqq{} \anchor{LANA}{2815--2819}.

Der Quotient $R^{\mathrm{ss}}$ des Berichts identifiziert adelische messbare
Mengen, wenn ihre Volumina gleich sind. Diese Definition beweist keine
Gleichheit der beiden konkreten Pilotklassen oder der Abbildungen in (9-1).
Die Gleichung bleibt im zitierten Bericht ein offenes Ziel (PDF S.~46, 48--49).
Es folgt eine vorgeschlagene objektweise Rollenzuordnung
(Tabelle~\ref{tab:lana_mapping}); sie liefert die fehlende getypte
Implikation zu Pflicht~7 nicht.

\begin{table}[htbp]
\centering
\small
\renewcommand{\arraystretch}{1.12}
\begin{tabular}{@{}P{0.27\textwidth}P{0.37\textwidth}P{0.23\textwidth}@{}}
\toprule
Pflicht & Objekt im Bericht & Anker \\
\midrule
3 --- Container, typisierte Einbettungen & das gemeinsame Ziel mit seinen zwei
  typisierten Abbildungen & 2790--2812 \\
4 --- erlaubte Identifikation & stipulierte Quellidentifikationen plus die Wahl
  von $S$ & 2937--2939 \\
5 --- getrenntes Budget & zwei Indeterminacies auf der anabelischen Seite, die
  dritte als Ambiguität der $S$-Wahl & 2786, 2973--2974 \\
\textbf{6 --- kompatible Auswertung} & das gemeinsame Ziel ist
  konstruktionsgemäß der Volumen-Auswertungsquotient, aber \textbf{die verlangte
  Gleichheit der beiden Auswertungs-Isomorphismen --- eine
  Kompatibilitätsidentität im Sinne dieser Pflicht --- ist gefordert und
  unbewiesen} & 2790--2794, 2815--2819, 2988--2990 \\
7 --- Ordnungskompatibilität & hier keine getypte Implikation belegt;
  Objekt, normierte Auswertung und Halbgeradenschranke bleiben zu verbinden & 2815--2819 \\
\bottomrule
\end{tabular}
\caption{Abgleich zwischen den Objekten des Formalisierungsberichts (Projekt LANA) und den Zertifikatspflichten.}
\label{tab:lana_mapping}
\end{table}

Drei Aussagen über diese Konvergenz, und bewusst keine vierte.

\begin{enumerate}
\item \textbf{Vorgeschlagene Rollenzuordnung, keine Validierung.} Die
  Abbildungen und ihre verlangte Kompatibilität können eine Frage an Pflicht~6
  motivieren. Das identifiziert sie nicht formal mit dem Zertifikatsvertrag.
\item \textbf{Die getypte Brücke fehlt weiterhin.} Eine Implikation von (9-1)
  zu Pflicht~7 verlangt die begründete Wahl von $S$, die relevanten Objekte und
  Quotientenrepräsentanten, eine gemeinsame normierte Log-Auswertung sowie
  Orientierung und Obergrenze der Zielhalbgerade. Die Gleichheit zweier
  abstrakter Abbildungen allein liefert diese Voraussetzungen nicht.
\item \textbf{Hier wird kein Ordnungsschluss abgeleitet.} Wären diese
  Voraussetzungen und eine geeignete Vergleichsschranke bewiesen, könnte
  eine kompatible Auswertung ein Ordnungsargument tragen. Weder diese Arbeit
  noch der zitierte Zwischenbericht liefert die vollständige Kette.
  Die Zuordnung bleibt ein Forschungsauftrag, keine bewiesene einseitige
  Implikation und keine Widerlegung.
\end{enumerate}

Auch die eigene Position des Berichts ist genau festzuhalten, weil sie dem
Verdikt-Vokabular aus Abschnitt~\ref{sec:instrument} entspricht und keiner der
Parteipositionen: Die Kompatibilität sei \glqq not manifestly false\grqq, und einen Beweis
dafür gebe es nicht. Das ist dieselbe Unterscheidung, auf der das Ledger besteht.
\emph{Liefert nicht ist nicht widerlegt} --- und die Konvergenz, soweit sie
besteht, ist eine Konvergenz darauf, wo die Frage sitzt, nicht auf eine Antwort.

\paragraph{Was \glqq unabhängig\grqq{} hier heißt --- und was nicht.} Das Wort wird in einem
streng belegten Sinn verwendet. Das Zertifikat und seine neun Zeilen wurden im
Mai~2026 fixiert, bevor eine der vier hier berichteten Anwendungen durchgeführt
wurde; diese fanden im August~2026 statt; der Zwischenbericht erschien im Juli~2026
und zitiert diese Arbeit nicht. Das ist die gesamte Evidenz, und sie trägt genau
eine begrenzte chronologische Beobachtung --- dass das anfängliche Ledger dem
Bericht vorausgeht. Einfluss auf spätere Anwendungen oder Revisionen ist damit
nicht ausgeschlossen; auch die Abwesenheit von Kontakt oder Einfluss in einer
der Richtungen ist nicht belegt;
das wäre eine Aussage über Personen, die keine Lektüre der Dokumente entscheiden
könnte, und sie wird deshalb nicht getroffen.

Zwei weitere Einschränkungen gehören unmittelbar hierher. Der Bericht ist ein
Zwischenbericht, seine Autoren protokollieren internen Dissens, und unsere
Abbildung ist eine objektweise Zuordnung, kein Äquivalenzbeweis. Die Abbildung
setzt außerdem die Pflichten~1 und~2 voraus, statt sie auszuweisen: Die beiden
Objekte und die Zielhalbgerade gelten als bereits typisiert, und in der Tabelle
oben werden nur die Pflichten~3 bis~7 zugeordnet.

\subsection{Was nicht gezeigt ist}

\begin{enumerate}
\item \textbf{Fünf von neun Zeilen sind ungefüllt.} Die Zeilen R0, R1, R3, R5 und
  R7 waren in keiner der vier Anwendungen Gegenstand und sind in allen offen. R5
  wurde nur berührt, soweit R4 es erzwang. Ein Ledger mit vier gefüllten Zeilen
  ist ein Teil-Ledger, und die obigen Übersichtstabellen sind so zu lesen. Die
  vier sind zudem nicht von einheitlichem Status, und der Nenner sollte das nicht
  verdecken: Drei von ihnen --- R2, R4 und R6, letztere an beiden Hälften ihres
  Splits berichtet --- sind Kernzeilen, während R8 die getriggerte Erweiterung
  aus Abschnitt~\ref{sec:ledger} ist, die auf keine Pflicht abbildet und ohne
  ihren Auslöser außerhalb des Geltungsbereichs liegt. Eine Vollständigkeitsaussage
  über das Zertifikat ist eine Aussage über die acht Kernzeilen R0--R7, von denen
  drei gefüllt wurden.
\item \textbf{Keine Aussage über Korrektheit.} In keinem der vier Fälle wurde eine
  Beweiszeile nachvollzogen oder eine Abschätzung nachgerechnet. Insbesondere
  stellt Fall B fest, \emph{dass} und wo gerechnet wird, nicht, dass es stimmt.
\item \textbf{Keine Aussage über Reparierbarkeit.} Nichts hier sagt, dass die
  nicht eingelöste Pflicht nicht eingelöst werden kann; eine konkrete Route ---
  ob irgendein Text einen Satz enthält, der die Ordnungsrelation aus dem
  deklarierten Eigenschaftspaar ableitet, statt sie formal zu ziehen --- wurde als
  verbleibender Hebel benannt und nicht verfolgt.
\item \textbf{Keine Aussage über die Vermutung.} Das Wort \glqq Lücke\grqq{} bedeutet in
  dieser Arbeit \glqq eine Pflicht eines benannten Zertifikats, die ein zitierter
  Wortlaut nicht einlöst\grqq. Es bedeutet nicht, dass die $abc$-Vermutung offen oder
  geschlossen ist, und es bedeutet nicht, dass eine hier geprüfte Theorie falsch
  ist.
\end{enumerate}

% =============================================================================
\section{Einschränkungen}
\label{sec:limitations}
% =============================================================================

Dieser Abschnitt sagt, was die Arbeit nicht tut, in den Begriffen der
zugrundeliegenden Prüfungen.

\begin{enumerate}
\item \textbf{Kein Beweis wurde nachvollzogen.} Kein Satz einer geprüften Quelle
  wurde Zeile für Zeile verifiziert, keine Abschätzung nachgerechnet, keine
  Konstante geprüft. Klassifiziert wurde die Struktur der Begründung.
\item \textbf{Vier von neun Zeilen, davon drei Kernzeilen.} Siehe
  Abschnitt~\ref{sec:convergence}. Die Zeilen-Splits aus
  Abschnitt~\ref{sec:splits} waren nicht Bestandteil des fixierten Zertifikats
  --- zwei entstanden aus dem Gebrauch, der dritte aus der gehärteten Lesart von
  Pflicht~6 --- und sie sind nur im Nachfolgeinstrument~\cite{GeigerBCM2026}
  normativ. Die Fallberichte bleiben deshalb überall dort bei den ursprünglichen
  Zeilen, wo der Split am Befund nichts ändert, damit keiner Anwendung
  zugeschrieben wird, sie habe eine Zeile gefüllt, die es zum Zeitpunkt ihrer
  Durchführung noch nicht gab. Wo der Split etwas ändert, nutzen sie ihn und sagen
  es: Fall~B berichtet an R4a und R4b, Fall~A an R6a und R6b (die Trennung von
  Fall~C wird im Fließtext gegen die zusammengelegte Zeile geführt). Version 0.1 hielt
  ohne Einschränkung fest, die Fallberichte berichteten gegen die ursprünglichen
  Zeilen; das beschrieb die eigene Praxis dieser Arbeit zu weit, denn Fall~B zog
  seine einzige Verdikts-Änderung bereits aus dem R4-Split.
\item \textbf{Ein Verdikt ruht auf einer Rekonstruktion.} In Fall A hat die
  tragende Zeile bei der Extraktion ein Relationssymbol verloren. Die
  Rekonstruktion und ihre Gegenproben sind in
  Abschnitt~\ref{sec:reconstruction-A} offengelegt; die Allgemeinformel \glqq kein
  Schluss hängt an der Stellung eines Symbols\grqq{} wäre für diese Arbeit falsch und
  wird bewusst nicht verwendet.
\item \textbf{Fall C trägt zwei spätere Korrekturen in entgegengesetzte
  Richtungen} (Abschnitt~\ref{sec:caseC-nachtraege}), und der Fall ist mit beiden
  zu lesen. Die der zweiten zugrundeliegende Substanzprüfung ist ihrerseits eine
  Form- und Substanzprüfung, keine Widerlegung; sie ist als
  Teil~5~\cite{GeigerSubstance2026} erschienen, ebenso wie die der ersten
  zugrundeliegende Quellenprüfung~\cite{GeigerGNC2026}.
\item \textbf{Fall D beschreibt eine Fassung eines Dokuments.} Die zweite Arbeit
  erklärt sich selbst für überholt \anchor{Zhou-II}{10}. Ihre Befunde sind an die
  Fassung gebunden.
\item \textbf{Die Konvergenz aus Abschnitt~\ref{sec:convergence} ist
  objektweise}, nicht formal. Eine Äquivalenz zwischen der siebten Pflicht des
  Zertifikats und der Kompatibilitätsaussage des Berichts wird weder behauptet
  noch bewiesen, und die Implikation, die wir angeben, läuft nur in eine Richtung.
\item \textbf{Mathematische Sachfragen bleiben planmäßig offen}, nicht aus
  Versehen --- siehe Abschnitt~\ref{sec:cannot}, Punkt~2.
\item \textbf{Das Instrument ist über diese Anwendungen hinaus nicht validiert.}
  Vier Anwendungen in zwei Theorien sind eine kleine Stichprobe. Ob das Zertifikat
  in anderen Feldern nützlich unterscheidet, ist ungeprüft; das Rezept aus
  Abschnitt~\ref{sec:recipe} wird angeboten, damit es geprüft werden kann.
\item \label{lim:biblio} \textbf{Bibliographische Disziplin.} Jeder Eintrag wurde
  vor dieser Fassung gegen die Originalquelle geprüft (Verifikationslauf vom
  13.~August~2026). Die Titel von~\cite{ZhouI} und~\cite{ZhouII} waren zunächst
  weggelassen --- das dieser Arbeit zugrundeliegende Material protokollierte ihre
  Kennungen, ihren Autor und ihre Zugehörigkeit, zitierte ihre Titel aber nicht,
  und einen plausiblen Titel zu erfinden wäre schlechter als ihn wegzulassen ---
  und wurden anschließend aus den arXiv-Einträgen selbst ergänzt. Die
  gelegentlich neben~\cite{LANA2026} genannte Formalisierungsnotiz vom April 2026
  wird bewusst \emph{nicht} zitiert: Sie wurde gesichert, aber nicht gelesen, und
  eine ungelesene Quelle hat in einem Literaturverzeichnis nichts zu suchen.
\item \textbf{Textbasis sind Extrakte.} Alle Zitate stammen aus
  \texttt{pdftotext}-Extrakten; Zeilenangaben sind Auffindhilfen, keine
  Seitenzählung.
\end{enumerate}

% =============================================================================
\section{Fazit}
\label{sec:conclusion}
% =============================================================================

Die Meinungsverschiedenheit, von der diese Arbeit ausgeht, ist zunächst keine über
Mathematik. Sie handelt davon, was eine bestimmte Satzfolge \emph{ist}. Für diese
Art von Streit sind die üblichen Instrumente schlecht geeignet, weil sie eine
Frage stellen --- ist das korrekt? ---, die beide Seiten über verschiedene
Gegenstände beantworten, und weil sie dem Autoritätsargument eine Tür offenlassen,
die sich von innerhalb einer Prosabegutachtung nicht schließen lässt.

Das Brücken-Zertifikat ersetzt diese Frage durch eine, die ein Text über sich
selbst beantworten kann: Welche von sieben benannten Pflichten löst dieser
Wortlaut ein? Um diese Frage herum liegen Disziplinen, die die Antworten
nachprüfbar statt bloß zuversichtlich machen --- typisierte statt graduierter
Lücken, ein Import-Ledger, protokollierte Suchbegriffe für jeden Negativbefund,
offengelegte Rekonstruktionen und ein Verdikt-Vokabular, in dem \glqq nicht eingelöst\grqq
nicht stillschweigend zu \glqq widerlegt\grqq{} werden kann.

Viermal angewandt, hat das Instrument vier verschiedene Dinge getan, und das ist
das Ergebnis, das zu behalten sich lohnt. Es hat zwei nicht eingelöste Pflichten
an einem benannten Substep detektiert und sie voneinander getrennt --- eine
Kompatibilitätsidentität, die nicht ausgewiesen ist, und eine Ordnungsrelation,
die nicht aus ihr abgeleitet ist --- eine Unterscheidung, auf die es ankommt, weil
sie einen diffusen Verdacht in zwei bestimmte fehlende Aussagen verwandelt statt
in eine allgemeine Klage. Es hat den
einen Ort weggeschlossen, an den die Quelle selbst ihre Bezifferung auslagert,
indem es zeigte, dass die dortige Rechnung explizit, an den von ihr benannten
Knoten geschlossen und stromabwärts des strittigen Schritts liegt. Es hat
unterschieden, als es auf ein unabhängiges Vergleichsprogramm angewandt wurde, und
drei liefernde Zeilen sowie eine anders
geformte Lücke an anderer Adresse zurückgegeben --- und es hat, als eine genauere
Lesart derselben Quellen es verlangte, die eigene erste Codierung dieses Falls
revidiert, statt sie stehen zu lassen. Und es hat protokolliert, dass
zwei nachgelagerte Anwendungen die nicht eingelöste Pflicht in einem Satz erben,
ohne sie zu prüfen.

Was dieses Bild ändern würde, ist ein einzelnes, angebbares Objekt: ein Satz, der
die Ordnungsrelation aus der erlaubten Identifikation ableitet, statt sie formal
aus Eigenschaften zu ziehen, die einer Konstruktion zugesprochen werden. Dass ein
solcher Satz im geprüften Material nicht gefunden wurde, ist eine Aussage über
dieses Material. Ob er anderswo existiert oder bewiesen werden kann, ist eine
mathematische Frage --- und sie ist die Frage, die dieses Instrument in einer
Form übergeben soll, die präzise genug ist, um daran zu arbeiten.

% =============================================================================
% KI-OFFENLEGUNG -- Pipeline-Standard, eingebunden aus der kanonischen Vorlage
% (zweisprachig, bringt eigenes \section* mit). Die Vorlage nennt Rechenskripte
% und Läufe; diese Arbeit hat keine, weshalb der Einschränkungsabsatz darunter
% nach unten abweicht -- so verlangt es die Vorlage selbst. Die Einschränkung
% NICHT streichen: ohne sie behauptete die Offenlegung nicht getane Arbeit.
% =============================================================================
\phantomsection
\addcontentsline{toc}{section}{AI Disclosure / KI-Offenlegung}
% =============================================================================
% AI DISCLOSURE -- PIPELINE-STANDARD (seit 2026-08-12, User-Freigabe)
% =============================================================================
% Ersetzt die frueheren Stufen L1-L5+ (archiviert unter _archive/).
% EINE Fassung fuer alle Paper dieser Pipeline.
%
% ABWEICHEN NACH UNTEN: Wurde in einem Paper tatsaechlich deutlich weniger
% KI eingesetzt, wird NICHT auf eine andere Datei gewechselt, sondern die
% Formulierung hier abgeschwaecht -- z. B. "in erheblichem Masse in allen
% Phasen" ersetzen durch die tatsaechlich zutreffenden Phasen, und aus der
% Qualitaetssicherungs-Aufzaehlung streichen, was nicht stattfand.
% Der Grundsatz bleibt: ehrlich benennen, was war -- nicht abgrenzen.
%
% Verwendung:
%   \input{../../_templates/disclosure/ai_disclosure_STANDARD}
% =============================================================================

\section*{AI Disclosure / KI-Offenlegung}

% --- German / Deutsch ---

\noindent
KI-Systeme wurden in erheblichem Ma\ss{}e in allen Phasen dieser Arbeit
eingesetzt: Konzeption und Forschungsdiskussion, Literatur- und
Quellenarbeit, mathematische Analyse und Beweisarbeit, Erstellung und
Pr\"ufung der Rechenskripte, Ausf\"uhrung und Auswertung der
Berechnungsl\"aufe, Text, \"Ubersetzung und Satz.

\medskip

\noindent
Qualit\"atssicherung: Einsatz mehrerer unabh\"angiger Modelle
verschiedener Anbieter mit Gegenpr\"ufungen und adversarialen KI-Reviews;
nachvollziehbare, versionierte Rechenskripte mit gespeicherten
Ergebnisdateien und Pr\"ufsummen; fortlaufende Beweis- und
Registerdokumente; strikte Claim-Hygiene (Behauptungen nur mit
dokumentiertem Beleg, Nachtr\"age statt stiller Korrekturen); Changelogs
\"uber alle Versionen.

\bigskip

\noindent\rule{\textwidth}{0.4pt}

\bigskip

% --- English ---

\noindent
AI systems were used extensively in all phases of this work: conception
and research discussion, literature and source work, mathematical analysis
and proof development, creation and checking of the computational scripts,
execution and evaluation of the computational runs, text, translation, and
typesetting.

\medskip

\noindent
Quality assurance: use of several independent models from different
providers with counter-checks and adversarial AI reviews; traceable,
versioned computational scripts with stored result files and checksums;
continuously maintained proof and register documents; strict claim hygiene
(assertions only with documented evidence, addenda instead of silent
corrections); changelogs across all versions.


\medskip

\noindent
\emph{Einschränkung für diese Arbeit:} Zu dieser Arbeit gehören keine
Rechenskripte und keine Berechnungsläufe; die entsprechenden Punkte des
Standardtextes (Erstellung und Prüfung der Rechenskripte, Ausführung und
Auswertung der Berechnungsläufe, versionierte Skripte mit Ergebnisdateien und
Prüfsummen) treffen hier nicht zu. Die Arbeit besteht aus Textbefunden an lokalen
PDF-Extrakten.

\medskip

\noindent
\emph{Restriction for this paper:} this work involved no computational scripts and
no computational runs; the corresponding items of the standard text (creation and
checking of computational scripts, execution and evaluation of computational runs,
versioned scripts with stored result files and checksums) do not apply here. The
work consists of textual findings against local PDF extracts.

% =============================================================================
% LITERATURVERZEICHNIS
% Die Provenienz der Metadaten steht pro Eintrag, in drei Graden: (i) wörtlich
% im dieser Arbeit zugrundeliegenden Material zitiert; (ii) Kennung und Version
% gegen den lokalen Extrakt-Kopf verifiziert; (iii) aus der quellengeprüften
% Bibliographie einer Schwesterarbeit desselben Programms übernommen. Kein
% Eintrag ohne einen dieser drei Grade.
% Pflicht-Quellencheck abgeschlossen am 2026-08-13 (WebFetch/WebSearch gegen
% arXiv, Project Euclid/EMS Press und die zitierten Manuskripte selbst) für
% alle 11 Einträge; siehe _results/QUELLENCHECK_BIBITEMS_2026-08-13.md. Die
% drei zuvor als vorläufig markierten Einträge (ZhouI, ZhouII, MFHMP) sind
% unten vervollständigt; alle übrigen Einträge wurden unverändert bestätigt.
% =============================================================================
\begin{thebibliography}{99}
\addcontentsline{toc}{section}{Literatur}
\small

% Provenienz (i): vollständige bibliographische Daten wörtlich im
% zugrundeliegenden Material zitiert, aus der Bibliographie von Zhou I
% (Extrakt Z. 2583-2584).
\bibitem{Mochizuki2021c}
S.~Mochizuki.
\newblock \emph{Inter-universal Teichmüller Theory III: Canonical Splittings of
the Log-theta-lattice}.
\newblock Publ. Res. Inst. Math. Sci. \textbf{57} (2021), Nr.~1/2, 403--626.
\newblock Hier zitiert als IUT-III; der verwendete Extrakt stammt von der
Institutsfassung 2020.

% Provenienz (ii)+(iii): Kennung und Rolle im zugrundeliegenden Material
% verifiziert; Zeitschriftendaten aus der quellengeprüften Bibliographie der
% Schwesterarbeit zur globalen Normierung im selben Programm.
\bibitem{Mochizuki2021d}
S.~Mochizuki.
\newblock \emph{Inter-universal Teichmüller Theory IV: Log-volume Computations
and Set-theoretic Foundations}.
\newblock Publ. Res. Inst. Math. Sci. \textbf{57} (2021), Nr.~1/2, 627--723.
\newblock Hier zitiert als IUT-IV; der verwendete Extrakt stammt von der
Institutsfassung 2020.

% Provenienz (iii). Die Bände I und II werden nur als Kontext zitiert, auf dem
% die geprüften Konstruktionen ruhen; kein Befund dieser Arbeit ruht auf ihnen.
\bibitem{Mochizuki2021a}
S.~Mochizuki.
\newblock \emph{Inter-universal Teichmüller Theory I: Construction of Hodge
Theaters}.
\newblock Publ. Res. Inst. Math. Sci. \textbf{57} (2021), Nr.~1/2, 3--207.

\bibitem{Mochizuki2021b}
S.~Mochizuki.
\newblock \emph{Inter-universal Teichmüller Theory II:
Hodge--Arakelov-theoretic Evaluation}.
\newblock Publ. Res. Inst. Math. Sci. \textbf{57} (2021), Nr.~1/2, 209--401.

% Provenienz (ii) für Kennung und Version (gegen den lokalen Extrakt-Kopf
% verifiziert), (iii) für den Titel.
\bibitem{JoshiATSIII}
K.~Joshi.
\newblock \emph{Construction of Arithmetic Teichmuller Spaces III: A `Rosetta
Stone' and a proof of Mochizuki's Corollary 3.12}.
\newblock arXiv:2401.13508, Version v4 (2025). Hier zitiert als ATS-III.

\bibitem{JoshiATSIV}
K.~Joshi.
\newblock \emph{Construction of Arithmetic Teichmuller Spaces IV: Proof of the
abc-conjecture}. Preliminary version for comments.
\newblock arXiv:2403.10430, Version v2 (2025). Hier zitiert als ATS-IV.

% Provenienz (i) für Kennung, Version, Autor und Zugehörigkeit, sämtlich im
% zugrundeliegenden Material zitiert; die Titel sind dort NICHT zitiert. Titel
% verifiziert am 2026-08-13 per WebFetch der arXiv-Abstract-Seiten
% (arxiv.org/abs/2503.14510 und arxiv.org/abs/2510.05448).
\bibitem{ZhouI}
Z.-P.~Zhou.
\newblock \emph{The inter-universal Teichmüller theory and new Diophantine
results over the rational numbers. I}.
\newblock Institute for Theoretical Sciences, Westlake University.
\newblock arXiv:2503.14510, Version v1 (2025). Hier zitiert als Zhou-I.

\bibitem{ZhouII}
Z.-P.~Zhou.
\newblock \emph{The inter-universal Teichmüller theory and new Diophantine
results over rational numbers. II}.
\newblock Institute for Theoretical Sciences, Westlake University.
\newblock arXiv:2510.05448, Version v1 (2025). Hier zitiert als Zhou-II. Das
Dokument bezeichnet sich selbst als ältere, noch zu aktualisierende Fassung.

% Provenienz (i): Titel und Spiegel im zugrundeliegenden Material protokolliert.
% Datum am 2026-08-13 von „17. Juli 2026" auf „16. Juli 2026" korrigiert, um mit
% der Fußzeile des Dokuments selbst übereinzustimmen, verifiziert per direktem
% WebFetch des PDFs
% (ncatlab.org/nlab/files/LANAProject-Report-July2026.pdf); zur Diskrepanz mit
% Sekundärberichten, die den 17. Juli nennen, siehe
% _results/QUELLENCHECK_BIBITEMS_2026-08-13.md.
\bibitem{LANA2026}
Project LANA.
\newblock \emph{Interim Report}.
\newblock 16.~Juli 2026 (Dokumentdatum; manche Sekundärberichte zur
Veröffentlichung nennen den 17.~Juli 2026). Hier zitiert als LANA.

% Provenienz (iii). Nur als Definition der Negativkontrolle zitiert; kein Befund
% dieser Arbeit verwendet die Quelle als Beleg. Kanonische Quelle und Datum
% verifiziert am 2026-08-13 per WebFetch des Manuskript-PDFs.
\bibitem{ScholzeStix2018}
P.~Scholze und J.~Stix.
\newblock \emph{Why abc is still a conjecture}.
\newblock Manuskript, 16.~Juli 2018. Verfügbar unter
\url{https://www.math.uni-bonn.de/people/scholze/WhyABCisStillaConjecture.pdf}.

% Provenienz (i): Titel und Autoren wörtlich im zugrundeliegenden Material
% zitiert, aus der Bibliographie von Zhou I (Extrakt Z. 2589-2591). Publikations-
% daten verifiziert am 2026-08-13 per WebFetch der Project-Euclid-Artikelseite
% (Kodai Mathematical Journal).
\bibitem{MFHMP}
S.~Mochizuki, I.~Fesenko, Y.~Hoshi, A.~Minamide und W.~Porowski.
\newblock \emph{Explicit estimates in inter-universal Teichmüller theory}.
\newblock Kodai Math. J. \textbf{45} (2022), Nr.~2, 175--236.
\newblock In Fall D als die Vergleichsarbeit für die nachgelagerten Konstanten
genannt.

% Eigene Vorarbeiten derselben Serie. Metadaten aus den Live-Zenodo-Records
% (geprüft 2026-08-13); es sind die beiden bereits erschienenen Teile der Serie.
\bibitem{GeigerGap2026}
L.~Geiger.
\newblock \emph{The Gap in IUTchIII, Corollary 3.12: From Attempted Repair to
Structural Diagnosis}.
\newblock Preprint, Version 6.3, Zenodo (2026).
\newblock \url{https://doi.org/10.5281/zenodo.21899747}.
\newblock Teil~1 der \emph{abc}-Gap-Serie; die veröffentlichte forensische
Untersuchung der in Fall~A geprüften Passage.

% Teil 2 derselben Serie, das Methodenpapier. Erschienen am 2026-08-16, während
% diese Version vorbereitet wurde; seine gehärtete sechste Vertragskomponente ist
% der Grund für die beiden in „Änderungen in Version 0.2" ausgewiesenen
% Umcodierungen.
\bibitem{GeigerBCM2026}
L.~Geiger.
\newblock \emph{The Bridge-Contract Criterion: A Form Audit for Transport
Arguments, with an Instantiation in the $abc$ Literature}.
\newblock Preprint, Version 0.2, Zenodo (2026).
\newblock \url{https://doi.org/10.5281/zenodo.21960477}.
\newblock Concept-DOI \url{https://doi.org/10.5281/zenodo.21960476}.
\newblock Teil~2 der \emph{abc}-Gap-Serie; das Methodenpapier, dessen gehärtete
sechste Vertragskomponente in Abschnitt~\ref{sec:obligations} dieser Version
übernommen wird.

\bibitem{GeigerSubstance2026}
L.~Geiger.
\newblock \emph{The Joshi ATS Series under the Bridge-Certificate Standard: A
Targeted Substance Audit and an All-Place Rigidity Classification of Effective
Normalization Weights}.
\newblock Preprint, Version 0.2, Zenodo (2026).
\newblock \url{https://doi.org/10.5281/zenodo.21956399}.
\newblock Concept-DOI \url{https://doi.org/10.5281/zenodo.21956398}.
\newblock Teil~5 derselben Serie; die in
Abschnitt~\ref{sec:caseC-nachtraege} zitierte Substanzprüfung.

\bibitem{GeigerGNC2026}
L.~Geiger.
\newblock \emph{Global Normalization in Joshi's Arithmetic Teichmüller Theory:
What the Identification Carries --- and What Carries the Theorem}.
\newblock Preprint, Version 1.1, Zenodo (2026).
\newblock \url{https://doi.org/10.5281/zenodo.21924254}.
\newblock Werkzeug-Supplement derselben Serie; die Quellenprüfung, die der ersten
der beiden Korrekturen in Abschnitt~\ref{sec:caseC-nachtraege} zugrunde liegt.

\end{thebibliography}

\end{document}